\documentclass[11pt,letterpaper]{article}
\usepackage[utf8]{inputenc}
\usepackage[T1]{fontenc}
\usepackage[spanish,es-noshorthands,es-tabla]{babel}
\usepackage{lmodern}
\usepackage[margin=2.2cm]{geometry}
\usepackage{amsmath,amssymb,amsthm,mathtools}
\usepackage{enumitem}
\usepackage{booktabs,array,tabularx}
\usepackage{multicol}
\usepackage{xcolor}
\usepackage[most]{tcolorbox}
\usepackage{tikz}
\usetikzlibrary{shapes.geometric,arrows.meta,positioning,calc}
\usepackage{fancyhdr}
\usepackage{titlesec}
\usepackage{listings}
\usepackage[hidelinks]{hyperref}

\definecolor{azul}{HTML}{1F4E79}
\definecolor{azulclaro}{HTML}{E8F0F8}
\definecolor{verde}{HTML}{2E7D32}
\definecolor{verdeclaro}{HTML}{EAF5EA}
\definecolor{naranja}{HTML}{B85C00}
\definecolor{naranjaclaro}{HTML}{FFF3E0}
\definecolor{gris}{HTML}{555555}
\definecolor{rojo}{HTML}{B71C1C}

\titleformat{\section}{\Large\bfseries\color{azul}}{\thesection}{0.6em}{}
\titleformat{\subsection}{\large\bfseries\color{azul}}{\thesubsection}{0.5em}{}
\titleformat{\subsubsection}{\normalsize\bfseries\color{gris}}{\thesubsubsection}{0.5em}{}

\setlength{\parindent}{0pt}
\setlength{\parskip}{4pt}
\setlist{itemsep=2pt,topsep=3pt}

\pagestyle{fancy}
\fancyhf{}
\renewcommand{\headrulewidth}{0.4pt}

% ---------- Entornos ----------
\newcounter{ej}
\newtcolorbox{ejercicio}[2][]{enhanced,breakable,colback=white,colframe=azul,
  boxrule=0.6pt,arc=2pt,left=6pt,right=6pt,top=4pt,bottom=4pt,
  title={\textbf{#2}},fonttitle=\small\bfseries,coltitle=white,colbacktitle=azul,#1}
\newtcolorbox{solucion}[2][]{enhanced,breakable,colback=verdeclaro,colframe=verde,
  boxrule=0.6pt,arc=2pt,left=6pt,right=6pt,top=4pt,bottom=4pt,
  title={\textbf{#2}},fonttitle=\small\bfseries,coltitle=white,colbacktitle=verde,#1}
\newtcolorbox{nota}[1][]{enhanced,breakable,colback=naranjaclaro,colframe=naranja,
  boxrule=0.5pt,arc=2pt,left=6pt,right=6pt,top=3pt,bottom=3pt,#1}
\newtcolorbox{formulario}[1][]{enhanced,breakable,colback=azulclaro,colframe=azul,
  boxrule=0.5pt,arc=2pt,left=6pt,right=6pt,top=3pt,bottom=3pt,#1}

\newcommand{\dif}[1]{%
  \ifcase#1\or\textcolor[HTML]{66BB6A}{$\bigstar$}\or\textcolor[HTML]{FFA726}{$\bigstar\bigstar$}\or\textcolor[HTML]{EF5350}{$\bigstar\bigstar\bigstar$}\fi}
\newcommand{\fuente}[1]{\hfill{\footnotesize\textcolor{white!80!azul}{[#1]}}}
\newcommand{\ptst}[1]{{\footnotesize\textcolor{white!80!azul}{(#1 pts)}}}
\newcommand{\pts}[1]{\textcolor{gris}{(#1 pts)}}

% ---------- Lógica ----------
\newcommand{\imp}{\Rightarrow}
\newcommand{\sii}{\Leftrightarrow}
\newcommand{\ent}{\models}
\newcommand{\der}{\vdash}
\newcommand{\p}[1]{\ensuremath{\textit{#1}}}

% ---------- Árboles de juego ----------
\tikzset{
  maxn/.style={draw,regular polygon,regular polygon sides=3,minimum size=0.75cm,inner sep=0pt,fill=azulclaro},
  minn/.style={draw,regular polygon,regular polygon sides=3,shape border rotate=180,minimum size=0.75cm,inner sep=0pt,fill=naranjaclaro},
  chn/.style={draw,circle,minimum size=0.6cm,inner sep=0pt,fill=verdeclaro},
  hoja/.style={draw,rectangle,minimum size=0.5cm,inner sep=1pt,font=\small},
  gnode/.style={draw,circle,minimum size=0.75cm,inner sep=0pt,font=\small\bfseries},
  garc/.style={-{Stealth[length=2mm]},thick},
}

\lstset{basicstyle=\ttfamily\small,columns=fullflexible,frame=single,framerule=0.3pt,rulecolor=\color{gris},
  literate={á}{{\'a}}1 {é}{{\'e}}1 {í}{{\'i}}1 {ó}{{\'o}}1 {ú}{{\'u}}1 {ñ}{{\~n}}1}

% ================= Figuras compartidas (guía y respuestas) =================

% ---- Gráfica G1: búsqueda no informada / informada ----
\newcommand{\figGUno}{%
\begin{tikzpicture}[scale=1.0]
  \node[gnode] (S) at (0,0) {S};
  \node[gnode] (A) at (2.4,1.3) {A};
  \node[gnode] (B) at (2.4,-1.3) {B};
  \node[gnode] (C) at (4.9,1.3) {C};
  \node[gnode] (D) at (4.9,-1.3) {D};
  \node[gnode,fill=verdeclaro] (G) at (7.3,0) {G};
  \draw[garc] (S) -- node[above left]{2} (A);
  \draw[garc] (S) -- node[below left]{1} (B);
  \draw[garc] (A) -- node[above]{2} (C);
  \draw[garc] (A) -- node[pos=0.3,right]{4} (D);
  \draw[garc] (B) -- node[below]{1} (D);
  \draw[garc] (B) to[bend right=50] node[below]{9} (G);
  \draw[garc] (C) -- node[above right]{4} (G);
  \draw[garc] (D) -- node[below right]{3} (G);
\end{tikzpicture}\\[2pt]
{\small\begin{tabular}{c|cccccc}$n$ & S & A & B & C & D & G\\\hline $h(n)$ & 4 & 2 & 3 & 1 & 2 & 0\end{tabular}}}

% ---- Gráfica G2: heurística admisible pero inconsistente ----
\newcommand{\figGDos}{%
\begin{tikzpicture}
  \node[gnode] (S) at (0,0) {S};
  \node[gnode] (A) at (2,1) {A};
  \node[gnode] (B) at (2,-1) {B};
  \node[gnode] (C) at (4,0) {C};
  \node[gnode,fill=verdeclaro] (G) at (6,0) {G};
  \draw[garc] (S) -- node[above left]{1} (A);
  \draw[garc] (S) -- node[below left]{1} (B);
  \draw[garc] (A) -- node[above right]{1} (C);
  \draw[garc] (B) -- node[below right]{2} (C);
  \draw[garc] (C) -- node[above]{3} (G);
\end{tikzpicture}\\[2pt]
{\small\begin{tabular}{c|ccccc}$n$ & S & A & B & C & G\\\hline $h(n)$ & 2 & 4 & 1 & 1 & 0\end{tabular}}}

% ---- Gráfica G3: examen simulado ----
\newcommand{\figGTres}{%
\begin{tikzpicture}
  \node[gnode] (S) at (0,0) {S};
  \node[gnode] (A) at (2.2,1.2) {A};
  \node[gnode] (B) at (2.2,-1.2) {B};
  \node[gnode] (C) at (4.6,0) {C};
  \node[gnode,fill=verdeclaro] (G) at (6.8,0) {G};
  \draw[garc] (S) -- node[above left]{1} (A);
  \draw[garc] (S) -- node[below left]{4} (B);
  \draw[garc] (A) -- node[right]{2} (B);
  \draw[garc] (A) -- node[above right]{5} (C);
  \draw[garc] (B) -- node[below right]{1} (C);
  \draw[garc] (C) -- node[above]{3} (G);
  \draw[garc] (B) to[bend right=35] node[below]{6} (G);
\end{tikzpicture}\\[2pt]
{\small\begin{tabular}{c|ccccc}$n$ & S & A & B & C & G\\\hline $h(n)$ & 6 & 5 & 2 & 3 & 0\end{tabular}}}

% ---- Gráfica bidireccional ----
\newcommand{\figBidir}{%
\begin{tikzpicture}
  \node[gnode] (S) at (0,0) {S};
  \node[gnode] (M) at (3.5,1.3) {M};
  \node[gnode] (A) at (2.2,-1.2) {A};
  \node[gnode] (B) at (4.8,-1.2) {B};
  \node[gnode,fill=verdeclaro] (G) at (7,0) {G};
  \draw[thick] (S) -- node[above left]{10} (M);
  \draw[thick] (M) -- node[above right]{10} (G);
  \draw[thick] (S) -- node[below left]{8} (A);
  \draw[thick] (A) -- node[below]{3} (B);
  \draw[thick] (B) -- node[below right]{8} (G);
\end{tikzpicture}}

% ---- Árbol T1 (MAX-MIN-MAX-hojas) ----
\newcommand{\figTUno}{%
\begin{tikzpicture}[level distance=1.25cm,
  level 1/.style={sibling distance=4.6cm},
  level 2/.style={sibling distance=2.3cm},
  level 3/.style={sibling distance=1.05cm}]
  \node[maxn] {R}
    child {node[minn] {A}
      child {node[maxn] {\scriptsize A1} child {node[hoja]{4}} child {node[hoja]{7}}}
      child {node[maxn] {\scriptsize A2} child {node[hoja]{2}} child {node[hoja]{9}}}}
    child {node[minn] {B}
      child {node[maxn] {\scriptsize B1} child {node[hoja]{6}} child {node[hoja]{1}}}
      child {node[maxn] {\scriptsize B2} child {node[hoja]{3}} child {node[hoja]{8}}}}
    child {node[minn] {C}
      child {node[maxn] {\scriptsize C1} child {node[hoja]{5}} child {node[hoja]{10}}}
      child {node[maxn] {\scriptsize C2} child {node[hoja]{12}} child {node[hoja]{2}}}};
\end{tikzpicture}}

% ---- Árbol T2 (AIMA Fig. 6.2) ----
\newcommand{\figTDos}{%
\begin{tikzpicture}[level distance=1.3cm,
  level 1/.style={sibling distance=3.6cm},
  level 2/.style={sibling distance=1.0cm}]
  \node[maxn] {R}
    child {node[minn] {A} child {node[hoja]{3}} child {node[hoja]{12}} child {node[hoja]{8}}}
    child {node[minn] {B} child {node[hoja]{2}} child {node[hoja]{4}} child {node[hoja]{6}}}
    child {node[minn] {C} child {node[hoja]{14}} child {node[hoja]{5}} child {node[hoja]{2}}};
\end{tikzpicture}}

% ---- Árbol T3 (profundidad 4) ----
\newcommand{\figTTres}{%
\begin{tikzpicture}[level distance=1.2cm,
  level 1/.style={sibling distance=7.6cm},
  level 2/.style={sibling distance=3.8cm},
  level 3/.style={sibling distance=1.9cm},
  level 4/.style={sibling distance=0.9cm}]
  \node[maxn] {R}
    child {node[minn] {\scriptsize 1}
      child {node[maxn] {\scriptsize 11}
        child {node[minn] {\tiny 111} child {node[hoja]{2}} child {node[hoja]{5}}}
        child {node[minn] {\tiny 112} child {node[hoja]{8}} child {node[hoja]{1}}}}
      child {node[maxn] {\scriptsize 12}
        child {node[minn] {\tiny 121} child {node[hoja]{4}} child {node[hoja]{6}}}
        child {node[minn] {\tiny 122} child {node[hoja]{7}} child {node[hoja]{3}}}}}
    child {node[minn] {\scriptsize 2}
      child {node[maxn] {\scriptsize 21}
        child {node[minn] {\tiny 211} child {node[hoja]{9}} child {node[hoja]{2}}}
        child {node[minn] {\tiny 212} child {node[hoja]{1}} child {node[hoja]{6}}}}
      child {node[maxn] {\scriptsize 22}
        child {node[minn] {\tiny 221} child {node[hoja]{5}} child {node[hoja]{8}}}
        child {node[minn] {\tiny 222} child {node[hoja]{3}} child {node[hoja]{4}}}}};
\end{tikzpicture}}

% ---- Árbol expectiminimax ----
\newcommand{\figExp}{%
\begin{tikzpicture}[level distance=1.3cm,
  level 1/.style={sibling distance=4.6cm},
  level 2/.style={sibling distance=2.3cm},
  level 3/.style={sibling distance=1.0cm}]
  \node[maxn] {R}
    child {node[chn] {$c_1$}
      child {node[minn]{} child {node[hoja]{3}} child {node[hoja]{9}} edge from parent node[left,font=\scriptsize]{0.5}}
      child {node[minn]{} child {node[hoja]{6}} child {node[hoja]{4}} edge from parent node[right,font=\scriptsize]{0.5}}}
    child {node[chn] {$c_2$}
      child {node[minn]{} child {node[hoja]{2}} child {node[hoja]{10}} edge from parent node[left,font=\scriptsize]{0.5}}
      child {node[minn]{} child {node[hoja]{20}} child {node[hoja]{15}} edge from parent node[right,font=\scriptsize]{0.5}}}
    child {node[chn] {$c_3$}
      child {node[minn]{} child {node[hoja]{7}} child {node[hoja]{8}} edge from parent node[left,font=\scriptsize]{0.3}}
      child {node[minn]{} child {node[hoja]{4}} child {node[hoja]{5}} edge from parent node[right,font=\scriptsize]{0.7}}};
\end{tikzpicture}}

% ---- Árbol de tres jugadores ----
\newcommand{\figMaxN}{%
\begin{tikzpicture}[level distance=1.3cm,
  level 1/.style={sibling distance=6.2cm},
  level 2/.style={sibling distance=3.1cm},
  level 3/.style={sibling distance=1.55cm},
  every node/.style={font=\small}]
  \node[draw,circle,fill=azulclaro] {1}
    child {node[draw,circle,fill=naranjaclaro] {2}
      child {node[draw,circle,fill=verdeclaro] {3} child {node{$(1,2,6)$}} child {node{$(4,2,3)$}}}
      child {node[draw,circle,fill=verdeclaro] {3} child {node{$(6,1,2)$}} child {node{$(7,4,1)$}}}}
    child {node[draw,circle,fill=naranjaclaro] {2}
      child {node[draw,circle,fill=verdeclaro] {3} child {node{$(5,1,1)$}} child {node{$(3,5,2)$}}}
      child {node[draw,circle,fill=verdeclaro] {3} child {node{$(7,7,1)$}} child {node{$(5,4,5)$}}}};
\end{tikzpicture}}

% ---- Paisaje 1D ----
\newcommand{\figPaisaje}{%
\begin{tikzpicture}[yscale=0.35,xscale=0.62]
  \foreach \x/\y in {0/3,1/5,2/6,3/5,4/4,5/6,6/8,7/9,8/7,9/4,10/2,11/5,12/7,13/6,14/3}{
    \fill[azul!60] (\x-0.35,0) rectangle (\x+0.35,\y);
    \node[font=\scriptsize] at (\x,\y+0.6) {\y};
    \node[font=\scriptsize,text=gris] at (\x,-0.7) {\x};
  }
  \draw (-0.6,0) -- (14.6,0);
  \node[font=\scriptsize,anchor=east] at (-0.7,-0.7) {estado};
  \node[font=\scriptsize,anchor=east] at (-0.7,4) {$f$};
\end{tikzpicture}}

% ---- Tablero 4 reinas [2,1,4,4] ----
\newcommand{\figReinas}[4]{%
\begin{tikzpicture}[scale=0.55]
  \foreach \i in {0,...,3}{\foreach \j in {0,...,3}{
    \pgfmathparse{mod(\i+\j,2)==0 ? "gris!25" : "white"}
    \fill[\pgfmathresult] (\i,\j) rectangle (\i+1,\j+1);}}
  \draw (0,0) rectangle (4,4);
  \foreach \c/\r in {0/#1,1/#2,2/#3,3/#4}{\node at (\c+0.5,\r-0.5) {\Large$\mathbf{Q}$};}
  \foreach \c in {1,...,4}{\node[font=\scriptsize] at (\c-0.5,-0.4) {\c};}
  \foreach \r in {1,...,4}{\node[font=\scriptsize] at (-0.4,\r-0.5) {\r};}
\end{tikzpicture}}

% ---- 8-puzzle ----
\newcommand{\puzzle}[9]{%
\begin{tikzpicture}[scale=0.5]
  \draw (0,0) grid (3,3);
  \node at (0.5,2.5){#1};\node at (1.5,2.5){#2};\node at (2.5,2.5){#3};
  \node at (0.5,1.5){#4};\node at (1.5,1.5){#5};\node at (2.5,1.5){#6};
  \node at (0.5,0.5){#7};\node at (1.5,0.5){#8};\node at (2.5,0.5){#9};
\end{tikzpicture}}

% ---- Gato ----
\newcommand{\gato}[9]{%
\begin{tikzpicture}[scale=0.45]
  \draw (1,0)--(1,3) (2,0)--(2,3) (0,1)--(3,1) (0,2)--(3,2);
  \node at (0.5,2.5){#1};\node at (1.5,2.5){#2};\node at (2.5,2.5){#3};
  \node at (0.5,1.5){#4};\node at (1.5,1.5){#5};\node at (2.5,1.5){#6};
  \node at (0.5,0.5){#7};\node at (1.5,0.5){#8};\node at (2.5,0.5){#9};
\end{tikzpicture}}

% ---- Mapa de Australia (grafo de restricciones) ----
\newcommand{\figAustralia}{%
\begin{tikzpicture}[scale=0.9]
  \node[gnode] (WA) at (0,0.5) {\scriptsize WA};
  \node[gnode] (NT) at (1.6,1.6) {\scriptsize NT};
  \node[gnode] (SA) at (2.0,0) {\scriptsize SA};
  \node[gnode] (Q) at (3.6,1.4) {\scriptsize Q};
  \node[gnode] (NSW) at (4.0,-0.2) {\scriptsize NSW};
  \node[gnode] (V) at (3.0,-1.5) {\scriptsize V};
  \node[gnode] (T) at (4.6,-2.2) {\scriptsize T};
  \draw[thick] (WA)--(NT) (WA)--(SA) (NT)--(SA) (NT)--(Q) (SA)--(Q) (SA)--(NSW) (SA)--(V) (Q)--(NSW) (NSW)--(V);
\end{tikzpicture}}

\fancyhead[L]{\small\textcolor{gris}{Inteligencia Artificial 2027-I · Guía para el Examen 1}}
\fancyhead[R]{\small\textcolor{verde}{Respuestas}}
\fancyfoot[C]{\small\thepage}

\begin{document}

\begin{titlepage}
\centering
\vspace*{1.5cm}
{\color{verde}\rule{\linewidth}{1.2pt}}\\[0.6cm]
{\Huge\bfseries\color{verde} Guía para el Examen 1}\\[0.3cm]
{\LARGE Respuestas y soluciones desarrolladas}\\[0.2cm]
{\large Inteligencia Artificial · PCIC, UNAM · 2027-I}\\[0.4cm]
{\color{verde}\rule{\linewidth}{1.2pt}}\\[1.2cm]
\begin{minipage}{0.84\linewidth}\large
Este documento acompaña a la \emph{Guía de estudio para el Examen 1}. La numeración coincide
(A1–A8, B1–B32, C1–C20, D1–D7, E1–E5). En los ejercicios de procedimiento se muestran las trazas
completas; en los de demostración se dan pruebas autocontenidas, como se pidió en las tareas.

\vspace{0.5cm}
Todas las trazas numéricas (búsquedas, AC-3, alfa-beta, tablas de verdad, conteos) se
verificaron con programas independientes.

\vspace{0.5cm}
\textbf{Consejo:} usa estas soluciones para \emph{corregir}, no para leer. Si una respuesta no
coincide con la tuya, identifica en qué paso se separan.
\end{minipage}
\vfill
\end{titlepage}

\tableofcontents
\newpage

% ============================================================
\section{Parte A — Conceptos}

\begin{solucion}{A1 · Agentes e introducción}
\textbf{1.} Un agente racional, para cada secuencia posible de percepciones, elige la acción que se
espera maximice su medida de desempeño, dada la evidencia de las percepciones y su conocimiento
previo. No es omnisciencia: el omnisciente conoce el resultado \emph{real} de sus acciones; el
racional maximiza el resultado \emph{esperado}. Ejemplo: cruzar la calle tras mirar a ambos lados
es racional aunque después caiga un objeto del cielo; el mal resultado no vuelve irracional la decisión.

\textbf{2.} Pensar como humano (modelado cognitivo), pensar racionalmente (leyes del pensamiento,
lógica), actuar como humano (prueba de Turing), actuar racionalmente (agente racional). El curso
adopta el \emph{agente racional}: es más general que las «leyes del pensamiento» (inferir
correctamente es solo una forma de actuar bien; a veces hay que actuar sin prueba, p.\,ej.\ un
reflejo), y es más apto para el desarrollo científico que imitar humanos, porque la racionalidad
está definida matemáticamente (maximizar utilidad esperada).

\textbf{3.} La prueba de Turing evalúa si un interrogador, conversando por texto, distingue a la
máquina de un humano: mide conducta (actuar como humano). Requiere lenguaje natural, representación
de conocimiento, razonamiento y aprendizaje. La habitación china (Searle) sostiene que manipular
símbolos según reglas (sintaxis) no produce comprensión (semántica): una persona que sigue un
manual para responder en chino no entiende chino; por tanto pasar la prueba no implica tener
mente. Va contra la \emph{IA fuerte} (las máquinas que actúan inteligentemente de verdad piensan y
tienen mente); la \emph{IA débil} solo afirma que pueden actuar \emph{como si} fueran inteligentes.

\textbf{4.} (a) \textbf{F}: la racionalidad es relativa a lo percibido; la aspiradora con sensor
local es racional. (b) \textbf{V}: si la acción correcta depende de percepciones pasadas que ya no
son visibles, un agente sin memoria no puede actuar bien (p.\,ej.\ detenerse cuando ambas celdas
estén limpias si solo percibe su celda). (c) \textbf{V}: un entorno donde todas las acciones
reciben el mismo desempeño (o solo existe una acción). (d) \textbf{F}: la función de agente
recibe la \emph{historia} completa de percepciones; el programa recibe la percepción actual.
(e) \textbf{F}: hay funciones no computables (p.\,ej.\ decidir la parada) o que exigen memoria
ilimitada. (f) \textbf{V}: si todas las acciones son igual de buenas, elegir al azar es óptimo.
(g) \textbf{F}: maximiza lo esperado; puede perder por mala suerte (cartas).

\textbf{5.} El agente reflejo elige con reglas condición–acción sobre la percepción actual (a lo
más, con un modelo del estado actual); no evalúa consecuencias. El que planea se pregunta «¿qué
pasaría si…?», necesita un modelo de cómo evoluciona el mundo y una meta o utilidad. Un reflejo
\emph{puede} ser racional cuando el entorno es totalmente observable y la mejor acción depende solo
de la percepción actual (la aspiradora de dos celdas de AIMA §2.2 lo es bajo sus supuestos).
\end{solucion}

\begin{solucion}{A2 · Búsqueda no informada}
\textbf{1.} Estado inicial, $\textsc{Acciones}(s)$, modelo de transición $\textsc{Resultado}(s,a)$,
prueba de meta y costo de acción $c(s,a,s')$. Solución: secuencia de acciones del inicial a una
meta. Óptima: la de menor costo de camino entre todas.

\textbf{2.} Un \emph{estado} es una configuración del mundo; un \emph{nodo} es una estructura del
árbol: estado, padre, acción, costo $g$ y profundidad. Varios nodos pueden contener el mismo estado
(caminos distintos). El grafo de estados tiene cada estado una vez; el árbol tiene un nodo por
\emph{camino}. Un ciclo $A\to B\to A$ genera caminos arbitrariamente largos: árbol infinito.

\textbf{3.} El algoritmo es el mismo; cambia la disciplina de la frontera: DFS pila (LIFO), BFS cola
(FIFO), UCS prioridad $g$, voraz prioridad $h$, A* prioridad $g+h$.

\textbf{4.} (a) \textbf{V} si todos los costos son iguales (UCS expande por profundidad). (b)
\textbf{V}: primero-el-mejor con $f(n)=-\text{profundidad}(n)$. (c) \textbf{V}: $h\equiv0$.
(d) \textbf{V}: BFS ignora costos; su completitud solo requiere $b$ finito (lo que se pierde es la
optimalidad). (e) \textbf{F}: DFS puede tener suerte y hallar la meta en su primera rama.

\textbf{5.} El primer camino a la meta que se \emph{genera} puede no ser el más barato: en Rumania,
de Sibiu se genera Bucarest vía Fagaras con $g=310$ antes de hallarla vía Pitesti con $g=278$. Al
\emph{extraer} con prioridad $g$ (costos $\ge0$) el $g$ del nodo extraído ya es mínimo.

\textbf{6.} La mayoría de los nodos está en el último nivel: $b^d\ge 1+b+\dots+b^{d-1}$. IDS genera
$d\,b+(d-1)b^2+\dots+1\cdot b^d=O(b^d)$. Con $b=10,d=5$: 123\,450 frente a 111\,110 de BFS ($\approx$11\% más).

\textbf{7.} UCS expande todo nodo con $g<C^*$. Cada acción cuesta $\ge\epsilon$, así que esos nodos
están a profundidad $\le\lfloor C^*/\epsilon\rfloor$; contando el nivel donde se genera la meta,
$O(b^{1+\lfloor C^*/\epsilon\rfloor})$. Si $\epsilon\to0$ la cota explota; con costos cero puede haber
caminos infinitos de costo 0 y UCS pierde la completitud.
\end{solucion}

\begin{solucion}{A3 · Heurísticas}
\textbf{1.} Admisible: $0\le h(n)\le h^*(n)$. Consistente: $h(n)\le c(n,a,n')+h(n')$ para todo
sucesor y $h(\text{meta})=0$. Consistente $\Rightarrow$ admisible (C2a), no al revés. Ejemplo: en
$G_2$, $h(A)=4=h^*(A)$ pero $h(A)>c(A,C)+h(C)=2$.

\textbf{2.} Un problema relajado tiene menos restricciones en las acciones (supergrafo de aristas);
su costo óptimo es una heurística admisible (y consistente) del original. En el 8-puzzle: «una ficha
se mueve de $A$ a $B$ si $A$ es adyacente a $B$ y $B$ está vacía». Quitando ambas condiciones
(teletransporte) se obtiene $h_1$; quitando solo «$B$ vacía» se obtiene $h_2$.

\textbf{3.} $h_2$ domina a $h_1$ si $h_2(n)\ge h_1(n)$ para todo $n$ (ambas admisibles). A* expande
seguro todo nodo con $g(n)+h(n)<C^*$; con $h_2$ ese conjunto está contenido en el de $h_1$, así que
expande menos (salvo empates). Conviene si calcular $h_2$ no es mucho más caro.

\textbf{4.} No es óptima (Arad$\to$Bucarest: 450 frente a 418). En árbol es incompleta (puede ciclar:
Iasi–Neamt). En grafo finito es completa.

\textbf{5.} Con $h\equiv0$, A* es UCS. Con $h=h^*$ (y desempates hacia mayor $g$) expande solo nodos de
un camino óptimo.

\textbf{6.} Si A* genera $N$ nodos y la solución está a profundidad $d$, $b^*$ satisface
$N+1=1+b^*+\dots+(b^*)^d$. Lo ideal es $b^*\approx1$.
\end{solucion}

\begin{solucion}{A4 · CSP}
\textbf{1.} $X=\{X_1,\dots,X_n\}$, dominios $D_i$, restricciones $C$ (alcance + relación).
Consistente: no viola restricciones; completa: asigna todas las variables; solución: consistente y
completa.

\textbf{2.} Unaria: $SA\ne\text{verde}$; binaria: $SA\ne WA$; global: $\textit{Alldiff}(F,T,U,W,R,O)$
o una fila de Sudoku; suave (preferencia/costo): «el profesor prefiere dar clase en la mañana».

\textbf{3.} Revisión hacia adelante: al asignar $X$, borra valores inconsistentes solo de los vecinos
no asignados de $X$ (un paso). Consistencia de arcos (AC-3/MAC): propaga en cadena; si el dominio
de $Y$ cambia, revisa los arcos que apuntan a $Y$. AC detecta más fallos: con $WA=R$, $Q=G$, FC deja
$NT=\{B\}$ y $SA=\{B\}$ adyacentes sin notar el conflicto; AC sí. AC cuesta más por nodo.

\textbf{4.} MRV elige la variable con menos valores legales (\emph{falla primero}: si el fallo es
inevitable, descubrirlo pronto poda más). Grado: desempata con la variable con más restricciones con
variables no asignadas. LCV elige el valor que quita menos opciones a los vecinos (\emph{falla al
último}). Asimetría: \emph{todas} las variables deben asignarse, así que empezar por la difícil no
cuesta nada; de los valores basta que \emph{uno} funcione, así que conviene probar primero el más
prometedor. (Si se buscaran todas las soluciones, el orden de valores no importaría.)

\textbf{5.} En el nivel $\ell$ se puede elegir cualquiera de las $n-\ell$ variables y cualquiera de
$d$ valores: $\prod_\ell (n-\ell)d=n!\,d^n$ hojas. Como las asignaciones conmutan, basta fijar una
variable por nivel: $d^n$ hojas.

\textbf{6.} (i) Algún dominio vacío: no hay solución. (ii) Todos los dominios con un valor: es
solución. (iii) Dominios con varios valores: puede haber 0, 1 o muchas soluciones; hay que buscar.
\end{solucion}

\begin{solucion}{A5 · Búsqueda local y metaheurísticas}
\textbf{1.} Trabaja con estados completos y guarda solo el estado actual (o unos pocos); el camino no
importa; memoria constante. Es incompleta: puede atorarse en óptimos locales/mesetas y no explora
sistemáticamente, así que no puede probar que no hay solución.

\textbf{2.} Máximo local: pico más alto que sus vecinos pero no global. Meseta: región plana
(máximo plano u hombro); el algoritmo vaga sin dirección. Cresta: secuencia de máximos locales no
conectados directamente; todo movimiento individual baja.

\textbf{3.} $T$ alta: $e^{\Delta E/T}\approx1$, acepta casi todo (caminata aleatoria, exploración).
$T$ baja: casi nunca acepta empeorar (hill climbing, explotación). Con enfriamiento suficientemente
lento converge al óptimo global con probabilidad $\to1$ (tiempo infinito).

\textbf{4.} La lista tabú prohíbe regresar a estados (o atributos) recientes; evita ciclos y permite
aceptar movimientos peores para salir de óptimos locales. Muy corta: ciclos. Muy larga: prohíbe
demasiado y bloquea buenos movimientos.

\textbf{5.} Representación (cromosoma, gen, alelo), población, aptitud, selección (proporcional o
torneo), cruza, mutación, reemplazo y criterio de paro. $P_m$ muy alta: búsqueda aleatoria que
destruye bloques buenos. $P_m=0$: no aparecen alelos nuevos; convergencia prematura.

\textbf{6.} Promediado sobre \emph{todas} las funciones objetivo posibles, todos los algoritmos de
optimización de caja negra rinden igual. Consecuencia: la ventaja de un algoritmo proviene de
explotar la estructura del problema; hay que elegir y ajustar el método por clase de problemas.
\end{solucion}

\begin{solucion}{A6 · Adversarios y utilidad}
\textbf{1.} $S_0$, $\textsc{ToMove}(s)$, $\textsc{Acciones}(s)$, $\textsc{Resultado}(s,a)$,
$\textsc{EsTerminal}(s)$, $\textsc{Utilidad}(s,p)$; una solución es una estrategia $S\to A$. Suma cero:
la suma de utilidades es constante (lo que gana uno lo pierde el otro). Información perfecta: el
estado es totalmente observable (ajedrez sí; póker no).

\textbf{2.} Tiempo $O(b^m)$, espacio $O(bm)$. En ajedrez $b\approx35$, $m\approx 80$–100: del orden de
$10^{123}$ nodos o más; imposible.

\textbf{3.} $\alpha$: mejor valor (mayor) que MAX ya tiene asegurado en el camino actual; $\beta$: mejor
(menor) para MIN. Se poda un subárbol cuando se demuestra que un jugador nunca permitiría llegar ahí;
la raíz no cambia, pero el valor devuelto por un nodo podado es solo una cota (p.\,ej.\ «$\le 2$»).

\textbf{4.} Mejor primero: $O(b^{m/2})$ (factor efectivo $\sqrt b$; duplica la profundidad alcanzable);
orden aleatorio: $\approx O(b^{3m/4})$; peor orden: sin poda, $O(b^m)$.

\textbf{5.} Ordenar los terminales igual que la utilidad; ser rápida; en no terminales, correlacionar
con la probabilidad real de ganar. Efecto horizonte: un evento dañino inevitable queda más allá de la
profundidad de búsqueda y maniobras dilatorias lo «esconden», dando una evaluación engañosa.

\textbf{6.} Cuando hay azar (dados) o el oponente no es óptimo y se modela probabilísticamente.
Minimax solo compara (max/min), así que cualquier transformación creciente de las hojas preserva la
decisión; expectimax promedia, y los promedios dependen de las diferencias de magnitud: solo es
invariante ante transformaciones afines positivas (C13).

\textbf{7.} Ordenabilidad, transitividad, continuidad, sustituibilidad, monotonicidad (y
descomponibilidad). Teorema de von Neumann–Morgenstern: existe $U$ con $U(A)\ge U(B)\iff A\succeq B$
y $U([p_1,S_1;\dots;p_n,S_n])=\sum_i p_iU(S_i)$. Principio: elegir la acción de máxima utilidad esperada.
\end{solucion}

\begin{solucion}{A7 · Lógica proposicional}
\textbf{1.} Sintaxis: qué oraciones están bien formadas. Semántica: valor de verdad de cada oración en
cada modelo. Teoría de pruebas: reglas sintácticas para derivar oraciones.

\textbf{2.} Modelo: asignación de verdad a cada símbolo. $KB\ent\alpha$: $\alpha$ es verdadera en todo
modelo de $KB$ ($M(KB)\subseteq M(\alpha)$). Válida: verdadera en todo modelo. Satisfacible: verdadera
en alguno. Equivalentes: mismos modelos.

\textbf{3.} $\alpha\ent\beta$ $\iff$ no hay modelo con $\alpha$ verdadera y $\beta$ falsa $\iff$
$\alpha\wedge\neg\beta$ no tiene modelos $\iff$ $\neg\alpha\vee\beta$ es verdadera en todo modelo $\iff$
$\alpha\imp\beta$ es válida.

\textbf{4.} $\textsc{Tell}(f)$: «ya lo sabía» si $M(KB)\subseteq M(f)$; «no lo creo» si
$M(KB)\cap M(f)=\emptyset$; «aprendí algo» si $\emptyset\subsetneq M(KB)\cap M(f)\subsetneq M(KB)$ (y
$M(KB)$ se reduce a esa intersección). $\textsc{Ask}(f)$: «sí» si $KB\ent f$, «no» si $KB\ent\neg f$,
«no sé» si es contingente.

\textbf{5.} Correcto: $\{f:KB\der f\}\subseteq\{f:KB\ent f\}$. Completo: $\supseteq$. Modus ponens es
correcto pero no completo: $\{A\vee B,\ A\imp C,\ B\imp C\}\ent C$ y MP no lo deriva. Con KB de
cláusulas definidas y consultas atómicas, el encadenamiento con MP es completo.

\textbf{6.} Horn: disyunción con a lo más un literal positivo. Definida: exactamente uno
($p_1\wedge\dots\wedge p_k\imp q$ o un hecho $q$). Meta: ninguno ($p_1\wedge\dots\wedge p_k\imp\p{False}$).
$A\vee B$ tiene dos positivos: no es Horn.

\textbf{7.} La resolución no genera todas las consecuencias (no deriva $A\vee\neg A$ de una KB vacía),
pero es \emph{completa por refutación}: si $KB\wedge\neg\alpha$ es insatisfacible, deriva la cláusula
vacía. Como $KB\ent\alpha$ sii $KB\wedge\neg\alpha$ es insatisfacible, se niega la consulta y se busca $\square$.
\end{solucion}

\begin{solucion}{A8 · Lógica de primer orden}
\textbf{1.} Objetos, relaciones (predicados), funciones, variables y cuantificadores. Términos:
constantes, variables, funciones aplicadas a términos. Fórmulas: atómicas (predicado sobre términos,
igualdad), compuestas con conectivos, cuantificadas.

\textbf{2.} $\forall x\,(\p{Est}(x)\wedge\p{Listo}(x))$ afirma que \emph{todo objeto} es estudiante y
listo (demasiado fuerte). $\exists x\,(\p{Est}(x)\imp\p{Listo}(x))$ es verdadera en cuanto exista
\emph{un} no estudiante (implicación vacuamente verdadera): demasiado débil.

\textbf{3.} Eliminar $\exists$ sustituyendo la variable por un término nuevo; preserva la
satisfacibilidad (no la equivalencia). El testigo puede depender de las variables universales
externas: en $\forall x\,\exists y\,\p{Madre}(y,x)$ cada $x$ tiene su madre, $y=M(x)$; una constante
diría que todos tienen la misma madre.

\textbf{4.} $\theta$ tal que $p\theta=q\theta$; el MGU es aquel del cual cualquier otro unificador es una
instancia (único salvo renombres). La verificación de ocurrencia impide unificar $x$ con un término
que contenga a $x$ (como $F(x)$), lo cual requeriría un término infinito y sería incorrecto.

\textbf{5.} Resolución SLD sobre cláusulas de Horn: encadenamiento hacia atrás en profundidad, reglas
de arriba abajo y submetas de izquierda a derecha (sin verificación de ocurrencia). Puede ciclar
porque DFS es incompleta en espacios infinitos: una regla recursiva por la izquierda se reinvoca a sí
misma antes de probar otras alternativas, aunque la respuesta sea derivable.
\end{solucion}

\newpage
% ============================================================
\section{Parte B — Procedimientos}

\begin{solucion}{B1 · PEAS y propiedades}
\renewcommand{\arraystretch}{1.15}
{\small
\begin{tabularx}{\linewidth}{>{\raggedright}p{2.1cm}>{\raggedright\arraybackslash}X}
\toprule
\textbf{Fútbol} & P: goles a favor/en contra, ganar. E: cancha, balón, compañeros, rivales. A: piernas
(correr, patear, pasar). S: cámaras, sensores de contacto, acelerómetros. \emph{Parcial, multiagente
(cooperativo y competitivo), estocástico, secuencial, dinámico, continuo.}\\
\textbf{Libros por internet} & P: conseguir los libros deseados al menor costo. E: sitios web,
vendedores. A: seguir ligas, llenar formularios, pagar. S: páginas web. \emph{Parcial, un agente
(o multi si hay otros compradores), casi determinista, secuencial, estático (semi: cambian precios),
discreto.}\\
\textbf{Partido de tenis} & P: ganar. E: cancha, pelota, rival. A: moverse, golpear. S: cámaras,
tacto. \emph{Parcial, multiagente, estocástico, secuencial, dinámico, continuo.}\\
\textbf{Tenis contra pared} & P: mejorar desempeño (golpes seguidos). E: pared, pelota, cancha. A y S
como arriba. \emph{Parcial, un agente, estocástico, secuencial (casi episódico por golpe),
dinámico, continuo.}\\
\textbf{Diagnóstico médico} & P: paciente sano, costos bajos. E: paciente, hospital, personal. A:
preguntas, pruebas, diagnósticos, tratamientos. S: síntomas y hallazgos capturados. \emph{Parcial,
un agente, estocástico, secuencial, dinámico, continuo.}\\
\textbf{Subasta} & P: obtener artículos valiosos al menor precio. E: casa de subastas, otros postores.
A: pujar. S: anuncios y pujas de otros. \emph{Totalmente observable (subasta abierta), multiagente
(competitivo), estratégico, secuencial, dinámico, discreto.}\\
\bottomrule
\end{tabularx}}

\smallskip
Varias clasificaciones son discutibles; en el examen lo importante es \emph{justificar}.
\end{solucion}

\begin{solucion}{B2 · Misioneros y caníbales}
\textbf{(a)} Estado $(m,c,b)$: misioneros y caníbales en la orilla inicial, $b=1$ si el bote está
ahí. Inicial $(3,3,1)$; meta $(0,0,0)$. Acciones: llevar $(i,j)\in\{(1,0),(2,0),(0,1),(0,2),(1,1)\}$.
Transición: si $b=1$, $(m-i,c-j,0)$; si $b=0$, $(m+i,c+j,1)$. Precondiciones: conteos en $[0,3]$ y
seguridad en ambas orillas: $(m=0\vee m\ge c)$ y $(3-m=0\vee 3-m\ge 3-c)$. Costo 1 por cruce.

\textbf{(b)} $4\cdot4\cdot2=32$ estados; los pares $(m,c)$ seguros son 10 ($m\in\{0,3\}$ con cualquier
$c$, más $(1,1)$ y $(2,2)$), así que hay 20 estados legales, de los cuales 16 son alcanzables desde el
inicial (p.\,ej.\ $(3,3,0)$ y $(0,0,1)$ no lo son).

\textbf{(c)} 11 cruces, por ejemplo:
$(3,3,1)\xrightarrow{CC}(3,1,0)\xrightarrow{C}(3,2,1)\xrightarrow{CC}(3,0,0)\xrightarrow{C}(3,1,1)
\xrightarrow{MM}(1,1,0)\xrightarrow{MC}(2,2,1)\xrightarrow{MM}(0,2,0)\xrightarrow{C}(0,3,1)
\xrightarrow{CC}(0,1,0)\xrightarrow{M}(1,1,1)\xrightarrow{MC}(0,0,0)$.
BFS, IDS o UCS garantizan la óptima (costos unitarios).

\textbf{(d)} Porque hay que hacer movimientos que parecen retroceder (regresar dos personas a la
orilla inicial en el paso 6); las heurísticas humanas de «progreso monótono» engañan.
\end{solucion}

\begin{solucion}{B3 · Tamaños de espacios}
\textbf{(a)} $9!/2=181\,440$ (el libro imprime 181\,400 por errata) y $16!/2\approx1.05\times10^{13}$.
Las permutaciones se dividen en dos clases de paridad; los movimientos preservan un invariante de
paridad, así que solo la mitad es alcanzable desde un estado dado.

\textbf{(b)} Colocar 8 reinas en casillas vacías en secuencia: $64\cdot63\cdots57\approx1.8\times10^{14}$
secuencias. La formulación por columnas, sin ataques, reduce a 2\,057 estados.

\textbf{(c)} Cota simple: $9!=362\,880$ secuencias de jugadas (el número real de partidas, que terminan
antes al ganar alguien, es 255\,168). Hay 5\,478 posiciones legales alcanzables (765 salvo simetrías).
\end{solucion}

\begin{solucion}{B4 · Trazas en $G_1$}
\textbf{(a) BFS.} S $\to$ fr.\,[A,B]; A $\to$ [B,C,D]; B $\to$ [C,D,G] (D ya estaba); C; D; G.
Orden: S, A, B, C, D, G. Camino \textbf{S–B–G}, costo \textbf{10} (menos acciones, no menor costo).

\textbf{(b) DFS.} S (apila B, A); A (apila D, C); C (apila G); G. Orden: S, A, C, G. Camino
\textbf{S–A–C–G}, costo \textbf{8}.

\textbf{(c) IDS.} $L=0$: S. $L=1$: S, A, B. $L=2$: S, A, C, D, B, D, G $\Rightarrow$ camino
\textbf{S–B–G}, costo \textbf{10}.

\textbf{(d) UCS.}
\begin{center}\small
\begin{tabular}{clll}
\toprule
Paso & Extrae ($g$) & Frontera después & Comentario\\
\midrule
1 & S (0) & A(2), B(1) & \\
2 & B (1) & A(2), D(2), G(10) & \\
3 & A (2) & D(2), C(4) & D vía A costaría 6: no mejora\\
4 & D (2) & C(4), G(5) & G mejora de 10 a 5\\
5 & C (4) & G(5) & G vía C = 8: no mejora\\
6 & G (5) & — & meta\\
\bottomrule
\end{tabular}
\end{center}
Camino \textbf{S–B–D–G}, costo \textbf{5} (óptimo).
\end{solucion}

\begin{solucion}{B5 · Voraz, A* y propiedades de $h$}
\textbf{(a)} $h^*$: S=5, A=6, B=4, C=4, D=3, G=0.

\textbf{(b)} Admisible: $4\le5$, $2\le6$, $3\le4$, $1\le4$, $2\le3$, $0\le0$. Consistente:
S$\to$A: $4\le2+2$; S$\to$B: $4\le1+3$; A$\to$C: $2\le2+1$; A$\to$D: $2\le4+2$; B$\to$D: $3\le1+2$;
B$\to$G: $3\le9$; C$\to$G: $1\le4$; D$\to$G: $2\le3$. \textbf{Sí es consistente} (con igualdad en tres arcos).

\textbf{(c) Voraz:} S(4) $\to$ A(2), B(3); extrae A $\to$ C(1), D(2); extrae C $\to$ G(0); extrae G.
Camino S–A–C–G, costo 8: \textbf{no óptimo}.

\textbf{A*:}
\begin{center}\small
\begin{tabular}{clll}
\toprule
Paso & Extrae ($g$, $f$) & Frontera después ($f$) & Comentario\\
\midrule
1 & S (0, 4) & A(4), B(4) & \\
2 & A (2, 4) & B(4), C(5), D(8) & empate A/B: alfabético\\
3 & B (1, 4) & D(4), C(5), G(10) & D mejora: $g=2$, $f=4$\\
4 & D (2, 4) & C(5), G(5) & G mejora: $g=5$\\
5 & C (4, 5) & G(5) & G vía C: $g=8$, no mejora\\
6 & G (5, 5) & — & meta\\
\bottomrule
\end{tabular}
\end{center}
Camino S–B–D–G, costo \textbf{5}: óptimo, como garantiza la consistencia. Observa que los valores
$f$ extraídos nunca decrecen: $4,4,4,4,5,5$.
\end{solucion}

\begin{solucion}{B6 · Rumania}
\textbf{(a)} Extracciones UCS: Sibiu 0, Rimnicu 80, Fagaras 99, Arad 140, Oradea 151, Pitesti 177,
Zerind 215, Craiova 226, Timisoara 258, \textbf{Bucarest 278}. Bucarest se genera primero al expandir
Fagaras con $g=310$; si se probara la meta al generar, se devolvería 310. Al expandir Pitesti se
encuentra $177+101=278<310$.

\textbf{(b) A* desde Lugoj:}
\begin{center}\small
\begin{tabular}{llll}
\toprule
Extrae & $g+h=f$ & Genera\\
\midrule
Lugoj & $0+244=244$ & Mehadia $70+241=311$, Timisoara $111+329=440$\\
Mehadia & $70+241=311$ & Drobeta $145+242=387$\\
Drobeta & $145+242=387$ & Craiova $265+160=425$\\
Craiova & $265+160=425$ & Pitesti $403+100=503$, Rimnicu $411+193=604$\\
Timisoara & $111+329=440$ & Arad $229+366=595$\\
Pitesti & $403+100=503$ & Bucarest $504+0=504$ (Rimnicu vía Pitesti 693: no mejora)\\
Bucarest & $504+0=504$ & meta\\
\bottomrule
\end{tabular}
\end{center}
Camino Lugoj–Mehadia–Drobeta–Craiova–Pitesti–Bucarest, costo \textbf{504}.

\textbf{(c)} Voraz desde Timisoara: Timisoara(329) $\to$ Lugoj(244) $\to$ Mehadia(241) $\to$ Drobeta(242)
$\to$ Craiova(160) $\to$ Pitesti(100) $\to$ Bucarest(0). Costo $111+70+75+120+138+101=\mathbf{615}$ frente
al óptimo 536 (Timisoara–Arad–Sibiu–Rimnicu–Pitesti–Bucarest).

\textbf{(d)} Con la distancia en línea recta \emph{a Fagaras}, Neamt parece más cercana que Vaslui, así que
se expande Neamt, que es un callejón sin salida; su único sucesor es Iasi, que vuelve a generar Neamt:
la búsqueda en árbol oscila Iasi–Neamt para siempre (incompleta). La ruta real pasa por Vaslui, que
«se aleja» primero.
\end{solucion}

\begin{solucion}{B7 · 8-puzzle, IDS y $b^*$}
\textbf{(a)} Inicio 1: todas las fichas fuera de lugar, $h_1=8$; $h_2=3+1+2+2+2+3+3+2=18$ (fichas 1 a 8).
Óptimo 26: $h_1$ captura el 31\% y $h_2$ el 69\% del costo real.

\textbf{(b)} Inicio 2: fuera de lugar 1–6 ($h_1=6$); Manhattan: $1+1+3+1+1+3+0+0=10$; costo óptimo \textbf{22}
(calculado por BFS). Nótese que un estado «casi ordenado» para el ojo humano está lejos de esta meta.

\textbf{(c)} BFS: $10+10^2+10^3+10^4+10^5=111\,110$. IDS: $5\cdot10+4\cdot10^2+3\cdot10^3+2\cdot10^4+10^5=123\,450$.

\textbf{(d)} $53=1+b^*+(b^*)^2+\dots+(b^*)^5$. Con $b^*=1.92$: $1+1.92+3.69+7.08+13.59+26.09\approx53.4$ \checkmark{}.
\end{solucion}

\begin{solucion}{B8 · Criptoaritmética}
Variables $F,T,U,W,R,O\in\{0,\dots,9\}$ y acarreos $C_1,C_2,C_3\in\{0,1\}$. Restricciones:
$\textit{Alldiff}(F,T,U,W,R,O)$; $O+O=R+10C_1$; $C_1+W+W=U+10C_2$; $C_2+T+T=O+10C_3$; $C_3=F$;
$T\ne0$, $F\ne0$.

Razonamiento: $F=C_3\le1$ y $F\ne0\Rightarrow F=1$, $C_3=1$, así que $2T+C_2\ge10\Rightarrow T\ge5$.
Probando $T=7$, $C_2=0$: $O=4$, luego $R=8$, $C_1=0$ y $U=2W$ con $W\notin\{1,4,7,8\}$ y $U$ distinto:
$W=2\Rightarrow U=4$ (choca con $O$); $W=3\Rightarrow U=6$ \checkmark{}. Solución: $734+734=1468$.
(Otra: $938+938=1876$.)
\end{solucion}

\begin{solucion}{B9 · AC-3 en Australia}
Dominios iniciales: $WA=\{G\}$, $V=\{R\}$, el resto $\{R,G,B\}$. Llamadas a \textsc{Revise} que borran:
\begin{enumerate}[nosep]
\item $NT\to WA$: borra $G$ $\Rightarrow NT=\{R,B\}$.
\item $SA\to WA$: borra $G$ $\Rightarrow SA=\{R,B\}$.
\item $SA\to V$: borra $R$ $\Rightarrow SA=\{B\}$.
\item $Q\to SA$: borra $B$ $\Rightarrow Q=\{R,G\}$.
\item $NSW\to SA$: borra $B$ $\Rightarrow NSW=\{R,G\}$.
\item $NSW\to V$: borra $R$ $\Rightarrow NSW=\{G\}$.
\item $NT\to SA$: borra $B$ $\Rightarrow NT=\{R\}$.
\item $Q\to NSW$: borra $G$ $\Rightarrow Q=\{R\}$.
\item $Q\to NT$: borra $R$ $\Rightarrow Q=\emptyset$: \textbf{inconsistencia detectada}.
\end{enumerate}
El orden exacto depende de la cola, pero el resultado (dominio vacío) no. Número de soluciones con 3
colores: $SA$ tiene 3 opciones; el anillo $WA$–$NT$–$Q$–$NSW$–$V$ es un camino que debe alternar los otros
2 colores (2 formas); $T$ es libre (3): $3\cdot2\cdot3=\mathbf{18}$.
\end{solucion}

\begin{solucion}{B10 · Revisión hacia adelante en 4 reinas}
\begin{enumerate}[nosep]
\item $Q_1=1$: $D_2=\{3,4\}$, $D_3=\{2,4\}$, $D_4=\{2,3\}$.
\item \quad $Q_2=3$: $D_3=\emptyset$ (3 está en diagonal, 2 y 4 también) $\Rightarrow$ fallo.
\item \quad $Q_2=4$: $D_3=\{2\}$, $D_4=\{3\}$.
\item \qquad $Q_3=2$: $D_4=\emptyset$ $\Rightarrow$ fallo. Sin más valores: retroceder hasta $Q_1$.
\item $Q_1=2$: $D_2=\{4\}$, $D_3=\{1,3\}$, $D_4=\{1,3,4\}$.
\item \quad $Q_2=4$: $D_3=\{1\}$, $D_4=\{1,3\}$.
\item \qquad $Q_3=1$: $D_4=\{3\}$.
\item \qquad\quad $Q_4=3$: \textbf{solución} $(2,4,1,3)$.
\end{enumerate}
\end{solucion}

\begin{solucion}{B11 · Ordenamiento}
\textbf{(a)} Todos los dominios tienen 3 valores: MRV empata. Por grado gana $SA$ (grado 5).
\textbf{(b)} $SA=\{B\}$ (vecino de $WA$ y $NT$), $Q=\{R,B\}$, los demás 3 valores: MRV elige $SA$.
\textbf{(c)} $Q\in\{R,B\}$. Con $Q=R$: $SA=\{B\}$, $NSW=\{G,B\}$ (3 valores restantes). Con $Q=B$: $SA=\emptyset$.
LCV elige \textbf{$Q=R$}.
\end{solucion}

\begin{solucion}{B12 · Hill climbing en 4 reinas}
\textbf{(a)} Pares que se atacan en $(2,1,4,4)$: (1,2) diagonal, (1,3) diagonal, (3,4) misma fila:
$h=3$.

\textbf{(b)} $h$ de los sucesores (fila nueva de la reina de cada columna):
\begin{center}\small
\begin{tabular}{c|cccc}
 & fila 1 & fila 2 & fila 3 & fila 4\\\hline
col.\,1 & 3 & * & \textbf{1} & 3\\
col.\,2 & * & 4 & 4 & 4\\
col.\,3 & 2 & 3 & 2 & *\\
col.\,4 & 3 & 3 & 4 & *
\end{tabular}
\end{center}
Mejor: mover la reina de la columna 1 a la fila 3: $(3,1,4,4)$ con $h=1$.

\textbf{(c)} En $(3,1,4,4)$ el único conflicto es (3,4); mover la reina de la columna 4 a la fila 2 da
$h=0$: $(3,1,4,2)$ es solución. Aquí hill climbing sí llega, en 2 pasos.
\end{solucion}

\begin{solucion}{B13 · Paisaje unidimensional}
\textbf{(a)} Desde 0: $0\to1\to2$ ($f=6$, máximo local). Desde 4: $4\to5\to6\to7$ ($f=9$, \textbf{óptimo
global}). Desde 10: $10\to11\to12$ ($f=7$, local). Desde 13: $13\to12$ ($f=7$).

\textbf{(b)} Tabú desde 2: vecinos 1 y 3 empatan con 5; desempate a la derecha: 3. Luego 2 es tabú,
se va a 4 (aunque empeora a 4); 3 tabú $\to$ 5; $\to6\to7\to8$. Recorrido $2,3,4,5,6,7,8$; mejor visto:
estado 7 con $f=9$. Hill climbing se habría detenido en 2.

\textbf{(c)} $\Delta E=7-9=-2$: $T=10$: $e^{-0.2}=0.819$; $T=1$: $e^{-2}=0.135$; $T=0.1$: $e^{-20}\approx2\times10^{-9}$.
\end{solucion}

\begin{solucion}{B14 · Genético en 8 reinas}
\textbf{(a)} \texttt{24748552}: pares atacados (1,8) misma fila, (2,4), (3,8), (6,7) diagonales: 4, así que
$28-4=24$. \texttt{32543213}: 17 pares atacados, $28-17=11$.

\textbf{(b)} Total $=78$: $24/78=30.8\%$, $23/78=29.5\%$, $20/78=25.6\%$, $11/78=14.1\%$.

\textbf{(c)} \texttt{327|52411} $\times$ \texttt{247|48552} $\Rightarrow$ \texttt{32748552} (aptitud 23) y
\texttt{24752411} (aptitud 22).

\textbf{(d)} El peor gana un torneo de 2 solo si ambos elegidos son él: $(1/4)^2=1/16=0.0625$.

\textbf{(e)} $(7/8)^8\approx0.344$.
\end{solucion}

\begin{solucion}{B15 · Minimax y alfa-beta en $T_1$}
\textbf{(a)} $A1=7$, $A2=9$, $A=7$; $B1=6$, $B2=8$, $B=6$; $C1=10$, $C2=12$, $C=10$; $R=10$: MAX juega hacia $C$.

\textbf{(b)}
\begin{itemize}[nosep]
\item $R(-\infty,\infty)\to A(-\infty,\infty)\to A1(-\infty,\infty)$: 4, 7 $\Rightarrow7$. En $A$: $\beta=7$.
\item $A2(-\infty,7)$: 2 ($v=2$, $\alpha=2$), 9 ($v=9\ge7$: corte, aunque ya no quedan hijos) $\Rightarrow A=7$. En $R$: $\alpha=7$.
\item $B(7,\infty)\to B1(7,\infty)$: 6, 1 $\Rightarrow 6$. En $B$: $v=6\le\alpha=7$ $\Rightarrow$ \textbf{se poda $B2$} (hojas 3 y 8).
\item $C(7,\infty)\to C1(7,\infty)$: 5, 10 $\Rightarrow 10$; en $C$: $\beta=10$. $C2(7,10)$: 12 $\ge10$ $\Rightarrow$ \textbf{se poda la hoja 2}.
\item $C=10$, $R=10$.
\end{itemize}
Hojas no evaluadas: \textbf{3, 8} (de $B2$) y \textbf{2} (de $C2$).
\end{solucion}

\begin{solucion}{B16 · Alfa-beta y orden}
\textbf{(a)} $A=3$, $B=2$, $C=2$, raíz 3. Tras $A$, $\alpha=3$. En $B$ la primera hoja da 2 $\le3$: se podan
\textbf{4 y 6}. En $C$ se ven 14, 5, 2 (el corte llega con la última). Se evalúan 7 de 9 hojas.

\textbf{(b)} Poner el mínimo de cada nodo MIN primero en $B$ y $C$: $C=(2,5,14)$ (y $B$ ya empieza con 2).
Se evalúan $3+1+1=\mathbf{5}$ hojas. Coincide con la cota de Knuth–Moore para orden perfecto:
$b^{\lceil m/2\rceil}+b^{\lfloor m/2\rfloor}-1=3+3-1=5$, es decir $O(b^{m/2})$. (Dentro del primer nodo MIN
no se puede ahorrar: con $\alpha=-\infty$ hay que ver todas sus hojas.)
\end{solucion}

\begin{solucion}{B17 · Expectiminimax}
\textbf{(a)} MIN: $\min(3,9)=3$, $\min(6,4)=4$, $\min(2,10)=2$, $\min(20,15)=15$, $\min(7,8)=7$, $\min(4,5)=4$.
$c_1=0.5\cdot3+0.5\cdot4=3.5$; $c_2=0.5\cdot2+0.5\cdot15=\mathbf{8.5}$; $c_3=0.3\cdot7+0.7\cdot4=4.9$. MAX elige $c_2$.

\textbf{(b)} Tratando el azar como MIN: $c_1=3$, $c_2=2$, $c_3=4$: elegiría $c_3$. La decisión cambia porque
el pesimista ignora que el resultado malo de $c_2$ es solo 50\% probable.

\textbf{(c)} Tras $c_1,c_2$, $\alpha=8.5$. En $c_3$, el primer nodo MIN vale 7; la cota superior es
$0.3\cdot7+0.7\cdot20=16.1>8.5$: no se poda. En el segundo nodo MIN, tras ver la hoja 4 sabemos que
vale $\le4$, así que $c_3\le0.3\cdot7+0.7\cdot4=4.9<8.5$: \textbf{se poda la hoja 5}.
\end{solucion}

\begin{solucion}{B18 · Tres jugadores}
Nodos del jugador 3 (maximiza $u_3$): $(1,2,6)$, $(6,1,2)$, $(3,5,2)$, $(5,4,5)$.
Nodos del jugador 2 (maximiza $u_2$): izquierda $(1,2,6)$ [2 > 1]; derecha $(3,5,2)$ [5 > 4].
Raíz (jugador 1): $(1,2,6)$ vs.\ $(3,5,2)$ $\Rightarrow$ \textbf{$(3,5,2)$}: obtienen 3, 5 y 2.

Alianzas: en juegos que no son de suma cero, dos jugadores pueden preferir el mismo resultado; p.\,ej.\
$(7,7,1)$ es mucho mejor para 1 y 2 que el resultado de equilibrio, pero el jugador 3 lo impide. Si 1 y 2
coordinaran (o si uno es fuerte, los débiles se unen contra él) aparecen alianzas, que pueden romperse
cuando dejan de convenir.
\end{solucion}

\begin{solucion}{B19 · Loterías}
\textbf{(a)} $EU(A)=0.8\sqrt4=1.6$; $EU(B)=\sqrt3\approx1.732$ $\Rightarrow B\succ A$. Equivalente cierto de $A$:
$\sqrt x=1.6\Rightarrow x=2.56$ (\$2\,560 $<$ valor esperado \$3\,200: aversión al riesgo).

\textbf{(b)} Con $U(\$0)=0$: $B\succ A\Rightarrow U(3k)>0.8\,U(4k)$. $C\succ D\Rightarrow0.2\,U(4k)>0.25\,U(3k)$;
multiplicando por 4, $0.8\,U(4k)>U(3k)$. Contradicción: ninguna $U$ explica ambas preferencias.

\textbf{(c)} Continuidad. $p\cdot10+(1-p)\cdot0=4\Rightarrow p=0.4$.
\end{solucion}

\begin{solucion}{B20 · Validez}
\begin{enumerate}[label=(\alph*),nosep]
\item \textbf{Válida}: es la contrapositiva, $(H\imp F)\equiv(\neg F\imp\neg H)$.
\item \textbf{Válida}: el antecedente es siempre falso.
\item \textbf{Insatisfacible}: $(A\imp B)\wedge(B\imp A)\equiv A\sii B$, conjuntada con su negación.
\item \textbf{Válida}: es la corrección de la regla de resolución.
\item \textbf{Válida} (ley de Peirce): si $A$ es falso, $A\imp B$ es verdadero, así que $(A\imp B)\imp A$ es falso y la implicación externa es verdadera; si $A$ es verdadero, el consecuente es verdadero.
\item \textbf{Contingente}: falsa con $A=F,B=T$; verdadera con $A=B=T$.
\item \textbf{Válida} (exportación): $\neg A\vee\neg B\vee C$ en ambos lados.
\item \textbf{Contingente}: falsa con $P=F,Q=T$ (es la «falacia de negar el antecedente»).
\end{enumerate}
\end{solucion}

\begin{solucion}{B21 · ¿Vinculación correcta?}
\begin{enumerate}[label=(\alph*),nosep]
\item \textbf{Sí}: \p{False} no tiene modelos. (b) \textbf{No}: \p{True} tiene modelos y \p{False} ninguno.
\setcounter{enumi}{2}
\item \textbf{Sí}: el único modelo ($A=B=T$) cumple $A\sii B$.
\item \textbf{No}: $A=B=F$ cumple $A\sii B$ y no $A\vee B$.
\item \textbf{Sí}: $A\sii B$ implica $A\imp B\equiv\neg A\vee B$.
\item \textbf{Sí}: de hecho son equivalentes: $(A\imp C)\vee(B\imp C)\equiv\neg A\vee\neg B\vee C$.
\item \textbf{Sí}: $(A\imp C)\wedge(B\imp C)\equiv(\neg A\vee C)\wedge(\neg B\vee C)\equiv C\vee(\neg A\wedge\neg B)$.
\item \textbf{Sí}: eliminación de la conjunción.
\item \textbf{No}: $A=T$, $C=F$, $D=T$, $E=F$ satisface la izquierda y falsifica $\neg D\vee E$.
\item \textbf{Sí}: $A=T,B=F$.
\item \textbf{Sí}: $A=B=T$.
\item \textbf{Sí}: para cada asignación de $A,B$ hay exactamente un valor de $C$ que satisface $(A\sii B)\sii C$; así que tiene la mitad de los modelos, igual que $A\sii B$ (4 de 8 sobre $A,B,C$).
\end{enumerate}
\end{solucion}

\begin{solucion}{B22 · Conteo de modelos}
(a) $B\vee C$ falla solo con $B=C=F$: $3\cdot4=\mathbf{12}$. (b) Falla solo con todo verdadero: $\mathbf{15}$.
(c) $A\wedge\neg B$ contradice $A\imp B$: $\mathbf{0}$. (d) $(A\wedge B)\vee(B\wedge C)\equiv B\wedge(A\vee C)$:
$1\cdot3\cdot2=\mathbf{6}$.
\end{solucion}

\begin{solucion}{B23 · Wumpus por revisión de modelos}
\textbf{(a)} $R_1:\neg P_{1,1}$; $R_2: B_{1,1}\sii(P_{1,2}\vee P_{2,1})$; $R_3: B_{2,1}\sii(P_{1,1}\vee P_{2,2}\vee P_{3,1})$;
$R_4:\neg B_{1,1}$; $R_5: B_{2,1}$. (Además $\neg P_{2,1}$, porque el agente estuvo ahí.)

\textbf{(b)} De $R_2,R_4$: $\neg P_{1,2}$. De $R_3,R_5,R_1$: $P_{2,2}\vee P_{3,1}$. Modelos $(P_{1,2},P_{2,2},P_{3,1})$:
$(F,T,F)$, $(F,F,T)$, $(F,T,T)$: \textbf{3 modelos} de 8.

\textbf{(c)} $\neg P_{1,2}$ se cumple en los 3: $KB\ent\neg P_{1,2}$ ([1,2] es segura).
$\neg P_{2,2}$ falla en 2 de ellos: \textbf{no} se sigue (tampoco su negación: es desconocido).
\end{solucion}

\begin{solucion}{B24 · CNF}
\textbf{(a)} $(B\imp(P_{12}\vee P_{21}))\wedge((P_{12}\vee P_{21})\imp B)$ $\to$
$(\neg B\vee P_{12}\vee P_{21})\wedge(\neg(P_{12}\vee P_{21})\vee B)$ $\to$
$(\neg B\vee P_{12}\vee P_{21})\wedge((\neg P_{12}\wedge\neg P_{21})\vee B)$ $\to$
$(\neg B_{11}\vee P_{12}\vee P_{21})\wedge(\neg P_{12}\vee B_{11})\wedge(\neg P_{21}\vee B_{11})$.

\textbf{(b)} $\neg(A\wedge B)\vee((\neg C\vee D)\wedge(\neg D\vee C))$ $\to$
$(\neg A\vee\neg B\vee\neg C\vee D)\wedge(\neg A\vee\neg B\vee\neg D\vee C)$.

\textbf{(c)} $\neg(\neg P\vee(Q\wedge R))\to P\wedge\neg(Q\wedge R)\to P\wedge(\neg Q\vee\neg R)$.

\textbf{(d)} $\neg V\vee(\neg N\vee R)$: una sola cláusula $\neg V\vee\neg N\vee R$.

\textbf{(e)} $\neg(\neg V\vee N)\vee R\to(V\wedge\neg N)\vee R\to(V\vee R)\wedge(\neg N\vee R)$.
No son equivalentes: con $V=F$, $N=F$, $R=F$, (d) es verdadera y (e) es falsa. La implicación no es asociativa.
\end{solucion}

\begin{solucion}{B25 · Resolución proposicional}
\textbf{(a)} Cláusulas: (1) $\neg B_{11}\vee P_{12}\vee P_{21}$, (2) $\neg P_{12}\vee B_{11}$, (3) $\neg P_{21}\vee B_{11}$,
(4) $\neg B_{11}$, (5) $P_{12}$ (consulta negada). (5)+(2): $B_{11}$; con (4): $\square$.

\textbf{(b)} (1) $\neg A\vee B\vee C$, (2) $A$, (3) $\neg B$, (4) $\neg C$. (1)+(2): $B\vee C$; +(3): $C$; +(4): $\square$.
La KB es insatisfacible (y por ello vincula cualquier cosa).

\textbf{(c)} Símbolos $L$ (llueve), $P$ (paraguas), $M$ (me mojo), $E$ (me enfermo). Cláusulas:
(1) $\neg L\vee P\vee M$, (2) $\neg M\vee E$, (3) $L$, (4) $\neg E$, (5) $\neg P$ (negación).
(2)+(4): $\neg M$; (1)+$\neg M$: $\neg L\vee P$; +(3): $P$; +(5): $\square$.

\textbf{(d)} Símbolos $My$ (mítico), $Im$, $Ma$ (mamífero), $Ho$ (cuerno), $Mg$ (mágico). CNF:
(1) $\neg My\vee Im$, (2) $My\vee\neg Im$, (3) $My\vee Ma$, (4) $\neg Im\vee Ho$, (5) $\neg Ma\vee Ho$, (6) $\neg Ho\vee Mg$.
\emph{Cuerno:} añadir $\neg Ho$: con (4) $\neg Im$; con (5) $\neg Ma$; $\neg Im$+(1): $\neg My$; $\neg My$+(3): $Ma$; $Ma+\neg Ma$: $\square$.
\emph{Mágico:} añadir $\neg Mg$: con (6) $\neg Ho$ y se repite lo anterior. \emph{Mítico:} no se sigue: el modelo
$My=F,Im=F,Ma=T,Ho=T,Mg=T$ satisface la KB (y $My=T,Im=T,Ma=F,Ho=T,Mg=T$ también, así que tampoco se sigue que no sea mítico).
\end{solucion}

\begin{solucion}{B26 · Encadenamiento}
\textbf{(a)} Contadores iniciales: $P\imp Q$:1, $L\wedge M\imp P$:2, $B\wedge L\imp M$:2, $A\wedge P\imp L$:2,
$A\wedge B\imp L$:2. Agenda $[A,B]$.
\begin{itemize}[nosep]
\item $A$: $A\wedge P\imp L\to1$, $A\wedge B\imp L\to1$.
\item $B$: $B\wedge L\imp M\to1$, $A\wedge B\imp L\to0$ $\Rightarrow$ agenda $+L$.
\item $L$: $L\wedge M\imp P\to1$, $B\wedge L\imp M\to0$ $\Rightarrow +M$.
\item $M$: $L\wedge M\imp P\to0$ $\Rightarrow +P$.
\item $P$: $P\imp Q\to0\Rightarrow+Q$; $A\wedge P\imp L\to0$ ($L$ ya inferido).
\item $Q$: es la consulta. $KB\der Q$.
\end{itemize}
\textbf{(b)} $Q\leftarrow P\leftarrow L\wedge M$. Para $L$: la regla $A\wedge P\imp L$ pide $P$, que ya está en la
pila de metas: \emph{ciclo}, se descarta; la regla $A\wedge B\imp L$ se cumple con los hechos. Para $M$:
$B\wedge L\imp M$ con $B$ hecho y $L$ ya probado. Luego $P$ y $Q$. Hay que detectar metas repetidas en la pila.
\end{solucion}

\begin{solucion}{B27 · DPLL}
Cláusulas: (1) $A\vee\neg B$, (2) $\neg A\vee C\vee D$, (3) $B\vee\neg C$, (4) $\neg D\vee\neg A$, (5) $E\vee A$,
(6) $E\vee\neg C$, (7) $\neg A\vee\neg B\vee\neg C$.
\begin{enumerate}[nosep]
\item No hay unitarias. \textbf{Símbolo puro:} $E$ solo aparece positivo $\Rightarrow E=T$; se satisfacen (5), (6).
\item Ningún otro símbolo es puro ni hay unitarias: \textbf{ramificar} $A=T$.
\begin{itemize}[nosep]
\item Quedan $C\vee D$, $B\vee\neg C$, $\neg D$, $\neg B\vee\neg C$. Unitaria $\neg D\Rightarrow D=F$.
\item Queda $C$: unitaria $\Rightarrow C=T$. Quedan $B$ y $\neg B$: $B=T$ viola (7). \textbf{Fallo.}
\end{itemize}
\item \textbf{Ramificar} $A=F$: quedan (1) $\neg B$ y (3) $B\vee\neg C$. $C$ es puro (solo negativo) $\Rightarrow C=F$;
luego $\neg B$: $B=F$. Todas las cláusulas satisfechas: \textbf{terminación temprana}.
\end{enumerate}
Modelo: $A=F$, $B=F$, $C=F$, $E=T$, $D$ arbitrario.
\end{solucion}

\begin{solucion}{B28 · Traducciones}
\begin{enumerate}[label=(\alph*),nosep]
\item $\neg\forall x\,[\p{Est}(x)\imp(\p{Toma}(x,\p{Frances},\p{P2001})\wedge\p{Toma}(x,\p{Griego},\p{P2001}))]$.
\item $\exists x\,[\p{Est}(x)\wedge\p{Toma}(x,\p{Griego},\p{P2001})\wedge\neg\p{Aprueba}(x,\p{Griego},\p{P2001})
\wedge\forall y\,((\p{Est}(y)\wedge\p{Toma}(y,\p{Griego},\p{P2001})\wedge\neg\p{Aprueba}(y,\p{Griego},\p{P2001}))\imp y=x)]$.
\item $\forall x\,\exists y\,\p{Ama}(x,y)$ \;y\; $\exists y\,\forall x\,\p{Ama}(x,y)$. No son equivalentes: la segunda implica la primera, no al revés (C19).
\item $\forall x\,\forall y\,[(\p{Persona}(x)\wedge\p{Poliza}(y)\wedge\p{Cara}(y))\imp\neg\p{Compra}(x,y)]$.
\item $\exists x\,[\p{Barbero}(x)\wedge\forall y\,((\p{Hombre}(y)\wedge\neg\p{Afeita}(y,y))\imp\p{Afeita}(x,y))]$.
\item $\forall x\,\p{Politico}(x)\imp\big[\exists y\,(\p{Persona}(y)\wedge\forall t\,(\p{Tiempo}(t)\imp\p{Engaña}(x,y,t)))\wedge
\forall y\,(\p{Persona}(y)\imp\exists t\,(\p{Tiempo}(t)\wedge\p{Engaña}(x,y,t)))\wedge
\neg\forall y\,\forall t\,((\p{Persona}(y)\wedge\p{Tiempo}(t))\imp\p{Engaña}(x,y,t))\big]$.
(«A todas algún tiempo» también admite la lectura $\exists t\,\forall y$; conviene declararla.)
\item $\forall x\,\forall y\,(\neg\p{Ama}(y,y)\imp\neg\p{Ama}(x,y))$.
\end{enumerate}
\end{solucion}

\begin{solucion}{B29 · CNF en primer orden}
\begin{enumerate}[label=(\alph*),nosep]
\item $\neg\p{Persona}(x)\vee\p{Madre}(M(x),x)$ \quad ($M$ de Skolem, depende de $x$).
\item $\p{Ama}(y,C)$ \quad ($C$ constante de Skolem: el $\exists$ no está dentro de ningún $\forall$).
\item $\forall x\,[\neg\exists y\,P(x,y)\vee Q(x)]\to\forall x\,[\forall y\,\neg P(x,y)\vee Q(x)]\to\neg P(x,y)\vee Q(x)$.
¡No hay Skolem!: el $\exists$ en el antecedente se vuelve $\forall$.
\item $\forall x\,[\exists y\,\neg P(x,y)\vee\exists z\,Q(x,z)]\to\neg P(x,F(x))\vee Q(x,G(x))$.
\item $\forall x\,[\neg\p{Est}(x)\vee\exists y\,(\p{Curso}(y)\wedge\neg\p{Aprueba}(x,y))]$ $\to$
$(\neg\p{Est}(x)\vee\p{Curso}(F(x)))\wedge(\neg\p{Est}(x)\vee\neg\p{Aprueba}(x,F(x)))$.
\item $(\p{Animal}(F(x))\vee\p{Ama}(G(x),x))\wedge(\neg\p{Ama}(x,F(x))\vee\p{Ama}(G(x),x))$.
\end{enumerate}
\end{solucion}

\begin{solucion}{B30 · Unificación}
\begin{enumerate}[label=(\alph*),nosep]
\item $\{y/\p{Juan},\ x/\p{Madre}(\p{Juan})\}$.
\item Falla: $x/\p{Juan}$ y luego $\p{Juan}$ vs.\ $\p{Isabel}$. Estandarizando $\p{Conoce}(x_2,\p{Isabel})$: $\{x_2/\p{Juan},\ x/\p{Isabel}\}$.
\item Falla: $x/F(y)$, luego $F(F(y))$ vs.\ $y$: la verificación de ocurrencia lo impide.
\item $\{x/F(A),\ y/A,\ z/A\}$.
\item $\{x/\p{Padre}(y),\ z/\p{Padre}(\p{Padre}(y))\}$.
\item $\{z/A,\ x/F(A),\ u/G(y)\}$.
\item Falla: $x/F(y)$ y luego $F(y)$ vs.\ $G(z)$ (símbolos de función distintos).
\item Falla: $x/G(z)$; $z/G(y)$ (así $x=G(G(y))$); después $G(G(y))$ vs.\ $A$.
\end{enumerate}
\end{solucion}

\begin{solucion}{B31 · Prolog}
\textbf{(a)} \texttt{r(Z)} $\to$ \texttt{p(Z), q(Z)}. \texttt{p(Z)} unifica con \texttt{p(a)}: \texttt{q(a)} falla
$\Rightarrow$ retroceso; \texttt{p(b)}: \texttt{q(b)} tiene éxito. Respuesta \texttt{Z = b}; al pedir más, falla.

\textbf{(b)} Regla 1: \texttt{W = maria}. Regla 2 con \texttt{Z = maria}: \texttt{ancestro(maria,W)} da
\texttt{W = luis} y, recursivamente, \texttt{W = ana}. Orden: \texttt{maria}, \texttt{luis}, \texttt{ana}; luego \texttt{false}.

\textbf{(c)} \texttt{ancestro(ana,W)}: la regla 1 falla (\texttt{ana} no es progenitora); la regla 2 llama
primero a \texttt{ancestro(ana,Z)}, idéntica a la meta original $\Rightarrow$ recursión infinita
(desbordamiento de pila). La semántica lógica es la misma; la procedural (DFS, izquierda a derecha) no.
\end{solucion}

\begin{solucion}{B32 · El coronel West}
\textbf{(a)} (1) $\p{American}(x)\wedge\p{Weapon}(y)\wedge\p{Sells}(x,y,z)\wedge\p{Hostile}(z)\imp\p{Criminal}(x)$;
(2) $\p{Owns}(\p{Nono},M_1)$, (3) $\p{Missile}(M_1)$ (instanciación existencial);
(4) $\p{Missile}(x)\wedge\p{Owns}(\p{Nono},x)\imp\p{Sells}(\p{West},x,\p{Nono})$;
(5) $\p{Missile}(x)\imp\p{Weapon}(x)$; (6) $\p{Enemy}(x,\p{America})\imp\p{Hostile}(x)$;
(7) $\p{American}(\p{West})$; (8) $\p{Enemy}(\p{Nono},\p{America})$.

\textbf{(b)} Iteración 1: (4) con $\{x/M_1\}$: $\p{Sells}(\p{West},M_1,\p{Nono})$; (5) con $\{x/M_1\}$: $\p{Weapon}(M_1)$;
(6) con $\{x/\p{Nono}\}$: $\p{Hostile}(\p{Nono})$. Iteración 2: (1) con
$\{x/\p{West},y/M_1,z/\p{Nono}\}$: $\p{Criminal}(\p{West})$. Punto fijo alcanzado.
\end{solucion}

\newpage
% ============================================================
\section{Parte C — Problemas tipo examen}

\begin{solucion}{C1 · Robot recolector}
\textbf{(a)} Sea $M$ el conjunto de casillas con moneda. Estado $s=(p,R)$ con $p$ la casilla del robot y
$R\subseteq M$ las monedas aún no recogidas. Inicial: $(H,\,M\setminus\{H\})$. Meta: $p=H$ y $R=\emptyset$.
Acciones: N, S, E, O si la casilla destino está en el tablero. Transición: $p'=p+\delta$,
$R'=R\setminus\{p'\}$. Costo 1.

\emph{Completitud.} Toda solución del problema original es una caminata $H=q_0,q_1,\dots,q_T=H$ por
casillas adyacentes que visita todas las de $M$. Le corresponde la sucesión de estados
$(q_t,R_t)$ con $R_t=M\setminus\{q_0,\dots,q_t\}$, que es un camino válido en el espacio (cada paso es
una acción legal) y termina en $(H,\emptyset)$ porque todas las monedas fueron visitadas. Recíprocamente,
un camino del inicial a la meta proyecta a una caminata que acaba en $H$ con $R_T=\emptyset$, es decir,
que visitó todas las monedas. Los costos coinciden (1 por paso), así que la correspondencia preserva la
optimalidad.

\textbf{(b)} Estados: $n^2\cdot2^k$. Factor de ramificación $b\le4$ (2 en esquinas, 3 en bordes).
Profundidad óptima: recoger las monedas en cualquier orden fijo con tramos de Manhattan
$\le2(n-1)$ cada uno y volver: $d\le 2(k+1)(n-1)$; también $d\le (n^2-1)+2(n-1)$ recorriendo todo el
tablero en serpiente y regresando. BFS en árbol genera $O(4^{d})$ nodos, $\le \sum_{i=0}^{d}4^i=(4^{d+1}-1)/3$;
en grafo, a lo más $n^2 2^k$ estados.

\textbf{(c)} BFS: memoria $O(b^d)$ (o $O(n^22^k)$ en grafo) y es óptima porque los costos son
uniformes. DFS: memoria $O(bm)$, pero en árbol hay ciclos (ir y volver), así que sin límite ni control de
ciclos puede no terminar; aun en grafo, devuelve el primer camino que encuentra: no es óptima.

\textbf{(d)} El estado $(p,R)$ se repite al llegar por secuencias distintas (NE vs.\ EN, ir y volver).
Búsqueda en grafo: conjunto de explorados sobre $(p,R)$; se descarta un estado ya explorado o en la
frontera con $g$ menor o igual. No se pierde la óptima: el costo futuro depende solo de $(p,R)$
(propiedad de Markov del estado); si un camino óptimo pasa por un estado repetido, sustituir su
prefijo por el primer camino a ese estado (de costo $\le$, porque BFS/UCS extrae en orden de $g$) da
otra solución de costo $\le$, es decir, también óptima.

\textbf{(e) Admisibilidad.} Toda solución desde $s$ debe entrar a cada $c\in R$ y terminar en $H$: la
parte de la caminata $p\leadsto c\leadsto H$ mide al menos $d(p,c)+d(c,H)$, pues Manhattan es la
distancia más corta en una cuadrícula sin obstáculos. Luego $h^*(s)\ge h_a(s)$ (y $\ge d(p,H)$ si
$R=\emptyset$). Para $h_b$: cada moneda exige entrar a su casilla, cada entrada es un movimiento distinto:
$h^*(s)\ge|R|$.

\emph{Consistencia de $h_a$.} Sea $s=(p,R)\to s'=(p',R')$ con costo 1, $|d(p,x)-d(p',x)|\le1$.
Para $c\in R'$: $d(p,c)+d(c,H)\le1+d(p',c)+d(c,H)\le1+h_a(s')$. Si se recogió $c'=p'$:
$d(p,c')+d(c',H)=1+d(p',H)\le1+h_a(s')$, porque $h_a(s')\ge d(p',H)$ (desigualdad del triángulo
$d(p',c)+d(c,H)\ge d(p',H)$, o definición si $R'=\emptyset$). Así $h_a(s)\le1+h_a(s')$; y $h_a=0$ en la meta.

\emph{Ninguna domina.} Con $H=(1,1)$, robot en $H$ y monedas en $(1,2),(2,1),(2,2)$: $h_a=4>3=h_b$.
Con $H$ en el centro de un tablero $3\times3$, robot en $H$ y monedas en los 4 vecinos: $h_a=2<4=h_b$.
Usaría $\max(h_a,h_b)$: admisible y consistente (C4), y domina a ambas.
\end{solucion}

\begin{solucion}{C2 · Consistencia, monotonía y no-reexpansión}
\textbf{(a)} Inducción sobre $k$, el número de acciones de un camino óptimo de $n$ a la meta. $k=0$: $n$
es meta, $h(n)=0=h^*(n)$. Paso: si el camino óptimo es $n\to n'\to\cdots$, el de $n'$ tiene $k-1$
acciones y es óptimo (subcamino de óptimo), así que $h(n')\le h^*(n')$ y
$h(n)\le c(n,a,n')+h(n')\le c(n,a,n')+h^*(n')=h^*(n)$. Si no hay camino, $h^*=\infty$. $\blacksquare$

\textbf{(b)} Si $n'$ es sucesor de $n$: $f(n')=g(n)+c(n,a,n')+h(n')\ge g(n)+h(n)=f(n)$. $\blacksquare$

\textbf{(c)} Por inducción fuerte sobre el orden de extracción: supongamos que todo nodo extraído antes
tuvo $g=g^*$, y que $n$ se extrae con $g(n)>g^*(n)$. Sea $P=(s=p_0,\dots,p_k=n)$ un camino óptimo a $n$
y $p_i$ el primer nodo de $P$ aún no expandido ($i\ge1$ porque $s$ ya se expandió). Su predecesor
$p_{i-1}$ se expandió con $g=g^*(p_{i-1})$, así que $p_i$ está en la frontera con
$g(p_i)\le g^*(p_{i-1})+c(p_{i-1},p_i)=g^*(p_i)$, es decir, $g(p_i)=g^*(p_i)$. Si $p_i=n$ contradice
$g(n)>g^*(n)$. Si no, aplicando la consistencia a lo largo de $P$ desde $p_i$ hasta $n$:
\[f(p_i)=g^*(p_i)+h(p_i)\le g^*(p_i)+c(P_{p_i\to n})+h(n)=g^*(n)+h(n)<g(n)+h(n)=f(n),\]
así que $p_i$ tendría que extraerse antes que $n$: contradicción. Por tanto todo nodo se extrae
con $g=g^*$: ningún camino posterior puede mejorarlo, no hace falta reabrirlo y \textbf{nunca se
reexpande}. En particular, la primera meta extraída tiene $g=g^*$; y como antes de extraerla hay en la
frontera un nodo del camino óptimo con $f\le C^*$ (admisibilidad), la meta extraída cumple
$g=f\le C^*$: es óptima. $\blacksquare$

\textbf{(d)} $S\xrightarrow{1}A\xrightarrow{1}G$ con $h(S)=2$, $h(A)=0$: admisible ($h^*(S)=2$,
$h^*(A)=1$) pero $h(S)=2>c(S,A)+h(A)=1$.
\end{solucion}

\begin{solucion}{C3 · Optimalidad de A* y consistencia}
\textbf{(a)} Sea $G^*$ una meta óptima ($g=C^*$) y $G_2$ una subóptima en la frontera ($g(G_2)>C^*$).
Mientras $G^*$ no se ha extraído, algún nodo $n$ de su camino óptimo está en la frontera (en árbol,
el sucesor de cada nodo expandido del camino se agrega). Con $h$ admisible:
$f(n)=g(n)+h(n)\le g(n)+h^*(n)=C^*$, y $f(G_2)=g(G_2)>C^*$. Así $f(n)<f(G_2)$: $n$ sale antes que $G_2$.
Aplicando el argumento a cada ancestro, todos los nodos del camino óptimo y $G^*$ salen antes que
$G_2$. $\blacksquare$

\textbf{(b)} $h^*$: S=5, A=4, B=5, C=3, G=0. Admisible. \textbf{Inconsistente}: $h(A)=4>c(A,C)+h(C)=2$.
A* en grafo:
\begin{center}\small
\begin{tabular}{clll}
\toprule
Paso & Extrae ($g$, $f$) & Frontera & Comentario\\
\midrule
1 & S (0, 2) & A(5), B(2) & \\
2 & B (1, 2) & C(4), A(5) & C con $g=3$\\
3 & C (3, 4) & A(5), G(6) & C se cierra con $g=3>g^*(C)=2$\\
4 & A (1, 5) & G(6) & C con $g=2$, pero ya está cerrado: se descarta\\
5 & G (6, 6) & — & \textbf{devuelve 6}\\
\bottomrule
\end{tabular}
\end{center}
El óptimo es S–A–C–G con costo \textbf{5}.

\textbf{(c)} En árbol, el segundo camino a C no se descarta: tras extraer A se agrega C$(g=2,f=3)$, que sale
de inmediato y genera G$(g=5,f=5)$, que sale antes que G$(6)$: devuelve \textbf{5}.

\textbf{(d)} (1) Usar una heurística consistente (p.\,ej.\ $h(A)\le2$, o aplicar \emph{pathmax}:
$f(n')\leftarrow\max(f(n'),f(n))$). (2) Permitir la reapertura: si se encuentra un camino más barato a
un nodo cerrado, regresarlo a la frontera (A* con reapertura es óptimo con $h$ admisible).
\end{solucion}

\begin{solucion}{C4 · Combinar heurísticas}
\textbf{(a) Verdadero.} $\max(h_1,h_2)(n)\le h^*(n)$ porque ambas lo son. Si ambas son consistentes y
$h=\max(h_1,h_2)$, sea $i$ con $h(n)=h_i(n)$: $h(n)=h_i(n)\le c+h_i(n')\le c+h(n')$.

\textbf{(b) Falso.} $S\xrightarrow{1}G$ con $h_1(S)=h_2(S)=1$: $h_1+h_2=2>1$.

\textbf{(c) Verdadero.} $\lambda h_1+(1-\lambda)h_2\le\lambda h^*+(1-\lambda)h^*=h^*$.

\textbf{(d) Falso} en general (mismo contraejemplo que (b)).

\textbf{(e) Falso.} $S\xrightarrow{1}A\xrightarrow{1}G$, $h_1=h^*$ ($2,1,0$, consistente) y $h_2=(2,0,0)$ admisible:
$\min=(2,0,0)$ y $2>1+0$. (El mínimo de dos \emph{consistentes} sí es consistente:
$h(n)\le h_j(n)\le c+h_j(n')=c+h(n')$ con $j$ tal que $h(n')=h_j(n')$.)
\end{solucion}

\begin{solucion}{C5 · Heurísticas del $N$-puzzle}
\textbf{(a)} Si un problema relajado tiene todas las acciones del original (y más), toda solución del
original lo es del relajado: el óptimo relajado es $\le h^*$. \emph{Relajación 1} (una ficha salta a
cualquier casilla): cada ficha mal colocada necesita exactamente un movimiento y los movimientos son
independientes: óptimo $=h_1$. \emph{Relajación 2} (una ficha se mueve a una casilla adyacente aunque
esté ocupada): las fichas se mueven independientemente y cada una necesita exactamente su distancia
Manhattan: óptimo $=h_2$.

\textbf{(b)} Cada acción desplaza una sola ficha una casilla horizontal o vertical: su distancia
Manhattan cambia exactamente en $\pm1$ y las demás no cambian. Luego $h_2(n)\le 1+h_2(n')=c+h_2(n')$ y
$h_2(\text{meta})=0$: consistente. Para $h_1$, una acción cambia el estado «bien/mal colocada» de a lo más
una ficha: $h_1(n)\le1+h_1(n')$: también consistente.

\textbf{(c)} Toda ficha mal colocada está a distancia Manhattan $\ge1$: $h_2\ge h_1$. Con $h$ consistente,
A* expande todo nodo con $f<C^*$ y ninguno con $f>C^*$. Si $n$ se expande con $h_2$ (y $f_2(n)<C^*$),
entonces $g(n)+h_1(n)\le g(n)+h_2(n)<C^*$: también se expande con $h_1$. Así, salvo nodos con $f=C^*$,
el conjunto expandido con $h_2$ está contenido en el de $h_1$.

\textbf{(d)} Si dos fichas están en su fila meta en orden invertido, en cualquier solución una de ellas
debe salir de la fila y volver: al menos 2 movimientos verticales que Manhattan no cuenta (su componente
vertical es 0). $lc$ cuenta el número mínimo de fichas a retirar para eliminar todos los conflictos de
cada fila/columna; cada retiro cuesta $\ge2$ movimientos extra en la dirección perpendicular, y los
extras por filas (verticales) y por columnas (horizontales) son movimientos distintos. Así
$h_3=h_2+2\,lc\le h^*$.
\end{solucion}

\begin{solucion}{C6 · A* ponderada}
Sea $P^*$ un camino óptimo de costo $C^*$. En búsqueda en árbol, mientras no se haya devuelto una meta,
algún nodo $n$ de $P^*$ está en la frontera con $g(n)=g^*(n)$. Como $w\ge1$ y $h$ admisible:
\[f_w(n)=g(n)+w\,h(n)\le w\,g(n)+w\,h^*(n)=w\,C^*.\]
Cuando se extrae una meta $G$: $h(G)=0$ y $g(G)=f_w(G)\le f_w(n)\le wC^*$ (si $G$ es el extremo de $P^*$, su costo es $C^*$). $\blacksquare$
$w=0$: UCS; $w=1$: A*; $w\to\infty$: búsqueda voraz.
\end{solucion}

\begin{solucion}{C7 · Búsqueda bidireccional}
\textbf{(a)} Adelante expande S: M(10), A(8). Atrás expande G: M(10), B(8). Adelante expande A: B($g_F=11$).
Atrás expande B: A($g_B=11$). Adelante expande M(10): G(20). Atrás expande M(10): M ya fue expandido
hacia adelante: \textbf{se detiene y devuelve S–M–G = 20}. Pero $C^*=8+3+8=\mathbf{19}$.

\textbf{(b)} Sea $P=(v_0=s,\dots,v_k=g)$ con $c(P)<\mu$, $a_i$ el costo del prefijo hasta $v_i$ y
$b_i=c(P)-a_i$. En Dijkstra, todo nodo con $d(s,v)<g^F_{\min}$ ya fue expandido hacia adelante (con su
$g_F$ exacto), y análogamente hacia atrás. Sea $i$ el mayor índice con $a_i<g^F_{\min}$ ($a_0=0$). Entonces
$v_i$ fue expandido hacia adelante y, por maximalidad, $a_{i+1}\ge g^F_{\min}$, de modo que
$b_{i+1}\le c(P)-g^F_{\min}<\mu-g^F_{\min}\le g^B_{\min}$: $v_{i+1}$ fue expandido hacia atrás con
$g_B(v_{i+1})\le b_{i+1}$. Al expandir $v_i$ se generó $v_{i+1}$ con $g_F\le a_{i+1}$; en el momento en
que $v_{i+1}$ tuvo ambos valores se actualizó $\mu\le a_{i+1}+b_{i+1}=c(P)<\mu$. Contradicción. Luego no hay
camino más barato que $\mu$: $\mu=C^*$. $\blacksquare$

\emph{Aplicación.} Al expandir A hacia adelante se genera B, que ya tiene $g_B=8$: $\mu=11+8=19$. Tras la
cuarta expansión, $g^F_{\min}=\min(M{:}10,B{:}11)=10$ y $g^B_{\min}=\min(M{:}10,A{:}11)=10$: $19\le20$,
\textbf{se detiene con 19} tras 4 expansiones.

\textbf{(c)} $f^F$ y $f^B$ estiman \emph{cada una} el costo total de la solución; sumarlas cuenta dos veces
el trayecto. Con heurísticas perfectas, $f^F_{\min}=f^B_{\min}=C^*$ y la condición $\mu\le2C^*$ aceptaría
cualquier primera solución, p.\,ej.\ 20 en vez de 19. Cotas válidas: $\max(f^F_{\min},f^B_{\min})$ y
$g^F_{\min}+g^B_{\min}+\epsilon$ ($\epsilon$ = costo mínimo de arista); MM usa el máximo de ellas.

\textbf{(d)} MM: $pr_F(n)=\max(f_F(n),2g_F(n))$, análogo atrás; expande el nodo de menor prioridad entre
ambas fronteras. Si $g_F(n)>C^*/2$ entonces $pr_F(n)>C^*$. Mientras no se haya certificado la solución, en
alguna frontera hay un nodo del camino óptimo con $g\le C^*/2$ en su dirección (el tramo no explorado del
camino óptimo tiene un extremo en cada mitad), con prioridad $\le\max(C^*,2\cdot C^*/2)=C^*$. Por tanto
MM nunca expande un nodo con $g>C^*/2$: las búsquedas se encuentran en la mitad (teorema de Holte et al., 2016).
\end{solucion}

\begin{solucion}{C8 · Verdadero o falso}
\begin{enumerate}[label=(\alph*),nosep]
\item \textbf{F}: la torre recorre varias casillas en un movimiento; Manhattan puede valer 14 cuando
bastan 2 movimientos. Admisible: 0, 1 (misma fila o columna) o 2.
\item \textbf{F}: se aplica sobre discretizaciones (rejillas, grafos de visibilidad, muestreo); es estándar en planificación de trayectorias.
\item \textbf{F}: sin consistencia puede devolver subóptimo ($G_2$, C3).
\item \textbf{F}: IDS devuelve la solución más somera, no la más barata.
\item \textbf{V}: si cada reinicio tiene probabilidad $p>0$ de éxito, $P(\text{fallar } r \text{ veces})=(1-p)^r\to0$.
\item \textbf{V}: número de intentos geométrico con media $1/p\approx7.1$ (unos 6 fallos y 1 éxito).
\item \textbf{F}: en el haz local los $k$ estados comparten información (se eligen los $k$ mejores de todos los sucesores), así que la búsqueda se concentra donde hay progreso.
\end{enumerate}
\end{solucion}

\begin{solucion}{C9 · Reyes en un tablero}
\textbf{(a-i)} Variables $K_1,\dots,K_k$; dominio: casillas $\{1..n\}^2$; para cada $i<j$:
$\max(|r_i-r_j|,|c_i-c_j|)\ge2$ (implica $K_i\ne K_j$). $\binom{k}{2}$ restricciones binarias.

\textbf{(a-ii)} $X_{r,c}\in\{0,1\}$; para cada par de casillas adyacentes (incluye diagonales)
$\neg(X_a\wedge X_b)$, y la restricción global $\sum X_{r,c}=k$. Pares adyacentes:
$2n(n-1)$ ortogonales $+\,2(n-1)^2$ diagonales $=2(n-1)(2n-1)$ restricciones binarias.

\textbf{(b)} En (i) los reyes son intercambiables: $k!$ soluciones equivalentes; se rompe la simetría
imponiendo $K_1<K_2<\dots<K_k$ (orden lexicográfico de casillas). Con $k$ pequeño, (i) con ruptura de
simetría tiene pocas variables; (ii) tiene $n^2$ variables pero restricciones locales que la revisión
hacia adelante propaga bien (poner un rey anula sus 8 vecinos). Ambas son defendibles; lo importante es
argumentar tamaño del dominio, densidad del grafo de restricciones y simetría.

\textbf{(c)} Estado: conjunto $S$ de casillas con rey. Acciones: agregar un rey en casilla vacía,
quitar un rey, mover un rey a otra casilla vacía. Resultado: el conjunto modificado. Objetivo (maximizar):
$f(S)=|S|-\lambda\cdot\text{conflictos}(S)$ con $\lambda>1$. Si hay un conflicto, quitar uno de los reyes da
$\Delta f\ge-1+\lambda>0$, así que todo máximo local está libre de ataques. Máximo teórico
$\lceil n/2\rceil^2$: se divide el tablero en $\lceil n/2\rceil^2$ bloques $2\times2$ (recortados en el borde),
cada uno admite a lo más un rey, y colocar reyes en las casillas (impar, impar) alcanza la cota.
\end{solucion}

\begin{solucion}{C10 · Límites de la consistencia y estructura}
\textbf{(a)} Triángulo $X,Y,Z$ con dominios $\{R,G\}$ y $\ne$ en cada arista: cada arco es consistente
(para $R$ existe $G$ y viceversa), pero un triángulo no se colorea con 2 colores.

\textbf{(b)} Hay $2c$ arcos. Un arco $(X_i,X_j)$ entra a la cola al inicio y cada vez que $D_j$ pierde un
valor: a lo más $d+1$ veces. \textsc{Revise} cuesta $O(d^2)$. Total $O(2c\cdot(d+1)\cdot d^2)=O(cd^3)$.

\textbf{(c)} Elegir una raíz y ordenar topológicamente $X_1,\dots,X_n$ (cada padre antes que sus hijos).
\emph{Hacia atrás:} para $j=n,\dots,2$, $\textsc{Revise}(\textit{Padre}(X_j),X_j)$, $O(d^2)$ cada una; si un dominio
queda vacío, no hay solución. \emph{Hacia adelante:} asignar a $X_1$ cualquier valor y a cada $X_i$ un valor
compatible con su padre, que existe por la consistencia dirigida. Sin retrocesos. Hay $n-1$ aristas:
$O(nd^2)$.

\textbf{(d)} En $n$-reinas las soluciones están densamente distribuidas: casi cualquier estado inicial
está cerca de una, y min-conflicts resuelve un millón de reinas en unos 50 pasos. En un Sudoku difícil hay
una sola solución y restricciones muy apretadas: el paisaje tiene muchos mínimos locales con pocos
conflictos, lejos de la solución; la propagación sistemática es mucho más eficaz.
\end{solucion}

\begin{solucion}{C11 · Minimax contra un MIN subóptimo}
\textbf{(a)} Sea $V(n)$ el valor minimax y $\sigma$ cualquier estrategia de MIN. Afirmación: si MAX juega
minimax, el resultado desde $n$ es $\ge V(n)$. Inducción en la altura. Hoja: igual. Nodo MAX: MAX elige
$c^*$ con $V(c^*)=V(n)$; por hipótesis el resultado es $\ge V(c^*)=V(n)$. Nodo MIN: MIN elige algún $c$; el
resultado es $\ge V(c)\ge\min_{c'}V(c')=V(n)$. $\blacksquare$

\textbf{(b)} Raíz MAX con dos jugadas. Izquierda: nodo MIN con hojas $(5,5)$, valor 5. Derecha: nodo MIN con
hojas $(0,10)$, valor 0. Minimax elige izquierda y obtiene 5. Si MIN siempre juega su acción de más a la
derecha (subóptimo), MAX obtiene 10 jugando a la derecha, que no es la jugada minimax.
\end{solucion}

\begin{solucion}{C12 · Corrección de alfa-beta}
\textbf{(a)} $n_1=\min(n_2,n_{21},\dots,n_{2b_2})$, $n_2=\max(n_3,n_{31},\dots,n_{3b_3})$, … Sustituyendo:
$n_1=\min(\dots,\max(\dots,\min(\dots,n_j,\dots),\dots),\dots)$.

\textbf{(b)} Sea $l_i$ el mejor valor (para el padre $n_{i-1}$) de los hermanos de $n_i$ ya evaluados a la
izquierda y $r_i$ el de los hermanos a la derecha. Entonces $n_{i-1}=\min(l_i,n_i,r_i)$ si $n_{i-1}$ es MIN y
$\max(l_i,n_i,r_i)$ si es MAX:
\[n_1=\min\big(l_2,r_2,\max\big(l_3,r_3,\min\big(l_4,r_4,\dots,\min(l_j,r_j,n_j)\dots\big)\big)\big).\]

\textbf{(c)} \emph{Lema.} Toda composición de funciones $y\mapsto\min(a,y)$ y $y\mapsto\max(b,y)$ es un
«recorte» $y\mapsto\min(H,\max(L,y))$ (constante si $L\ge H$), en particular monótona no decreciente.

Sea $\beta=\min\{l_i: n_{i-1}\text{ es MIN}\}$, alcanzado en el nivel $k$ ($l_k=\beta$). Escribamos
$n_1=F(\varphi_k(g(n_j)))$, donde $g$ agrupa los niveles entre $n_k$ y $n_j$, $\varphi_k(y)=\min(\beta,r_k,y)$
y $F$ los niveles superiores. Por el lema $g(x)=\min(H,\max(L,x))$. Para $x\ge\beta$ hay dos casos:
si $H>\beta$, entonces $g(x)\ge\min(H,x)\ge\beta$ y $\varphi_k(g(x))=\min(\beta,r_k)$; si $H\le\beta$, entonces
$g(x)=H$. En ambos casos $n_1$ es \textbf{constante} para $n_j\in[\beta,\infty)$: el valor exacto de $n_j$ no
importa en cuanto $n_j\ge\beta$. Así, $n_j$ solo puede afectar a $n_1$ si $n_j<\beta$.

Como $n_j$ es MAX, su valor es $\ge$ el de cualquier hijo ya evaluado; en cuanto un hijo da $v\ge\beta$,
sabemos que $n_j\ge\beta$ y los hijos restantes no pueden cambiar $n_1$: se podan. Es el \emph{corte
$\beta$}, y $\beta$ es exactamente el mejor valor de MIN a lo largo del camino.

\textbf{(d)} Si $n_j$ es MIN, se repite el argumento con $\alpha=\max\{l_i:n_{i-1}\text{ es MAX}\}$: $n_1$ es constante
para $n_j\in(-\infty,\alpha]$. Como $n_j$ solo puede bajar al ver más hijos, en cuanto un hijo da $v\le\alpha$ se
podan los restantes (\emph{corte $\alpha$}).
\end{solucion}

\begin{solucion}{C13 · Transformaciones de utilidad}
\textbf{(a)} Si $g$ es estrictamente creciente, $g(\max(a,b))=\max(g(a),g(b))$ y lo mismo con $\min$. Por
inducción desde las hojas, el valor de cada nodo del árbol transformado es $g$(valor original), y el
$\arg\max$ en la raíz no cambia porque $g$ preserva el orden. $\blacksquare$

\textbf{(b)} Nodo de azar: $\sum_ip_i(av_i+b)=a\sum_ip_iv_i+b$ (porque $\sum p_i=1$). Max/min se preservan
por (a) al ser $a>0$. Por inducción todos los valores se transforman afínmente y las decisiones no cambian. $\blacksquare$

\textbf{(c)} Raíz MAX con dos nodos de azar: $A=[0.5{:}\,1;\ 0.5{:}\,9]$ (media 5) y $B=[0.5{:}\,4;\ 0.5{:}\,5]$
(4.5): se elige $A$. Con $\sqrt{\cdot}$: $A=0.5\cdot1+0.5\cdot3=2$, $B=0.5\cdot2+0.5\cdot2.236=2.118$: se elige $B$.

\textbf{(d)} En un nodo de azar con hijos ya evaluados, cota superior
$\sum_{\text{vistos}}p_iv_i+(\sum_{\text{no vistos}}p_i)\,U$ y cota inferior análoga con $L$. Bajo un padre MAX con
valor $\alpha$: si la cota superior $\le\alpha$, se podan los hijos restantes; bajo un padre MIN con $\beta$: si la
cota inferior $\ge\beta$. Ejemplo en B17(c).
\end{solucion}

\begin{solucion}{C14 · Evaluación en el gato}
Casillas $(f,c)$ con $f,c\in\{0,1,2\}$; O en $(0,0)$, X en $(1,1)$.

\textbf{(a)} Líneas solo con X: fila 1, columna 1, antidiagonal ($X_1=3$). Solo con O: fila 0, columna 0
($O_1=2$). La diagonal principal tiene ambas. $\textsc{Eval}=3-2=\mathbf{1}$.

\textbf{(b)} La posición es simétrica respecto de la diagonal principal; jugadas distintas:
\begin{center}\small
\begin{tabular}{lcccc}
\toprule
Jugada de X & $X_2$ & $X_1$ & $O_1$ & \textsc{Eval}\\
\midrule
$(0,1)$ borde junto a O & 1 & 2 & 1 & 4\\
$(0,2)$ esquina libre & 1 & 3 & 1 & \textbf{5}\\
$(1,2)$ borde lejano & 1 & 3 & 2 & 4\\
$(2,2)$ esquina opuesta & 0 & 5 & 2 & 3\\
\bottomrule
\end{tabular}
\end{center}
A profundidad 1, X elige la esquina $(0,2)$.

\textbf{(c)} Mínimo sobre respuestas de O: $(0,1)$: respuestas valen $2,1,2,0,0,2$ $\Rightarrow$ \textbf{0} (O en
$(2,0)$ o $(2,1)$). $(0,2)$: $\Rightarrow-1$ (O bloquea en $(2,0)$: la antidiagonal queda cerrada y la columna 0 tiene dos O, $\textsc{Eval}=3-(3+1)=-1$). $(1,2)$: $\Rightarrow-1$ (O en $(1,0)$). $(2,2)$: $\Rightarrow-1$ (O en $(0,2)$ o $(2,0)$).
Decisión minimax a profundidad 2: \textbf{$(0,1)$} (o su simétrica $(1,0)$), con valor 0.

\textbf{(d)} No coinciden. A profundidad 1 la esquina crea una amenaza (dos en línea); a profundidad 2
se ve que O la bloquea y a la vez forma su propia línea doble. Es un ejemplo del efecto de la
profundidad (horizonte) sobre una evaluación imperfecta.
\end{solucion}

\begin{solucion}{C15 · Vinculación, resolución y Horn}
\textbf{(a)} ($\Rightarrow$) Si $\alpha\ent\beta$ y $m$ satisface $\alpha$, entonces satisface $\beta$, luego no satisface
$\neg\beta$: ningún modelo satisface $\alpha\wedge\neg\beta$. ($\Leftarrow$) Si $\alpha\wedge\neg\beta$ no tiene modelos, todo
modelo de $\alpha$ falsifica $\neg\beta$, es decir, satisface $\beta$. $\blacksquare$

\textbf{(b)} Sea $m$ modelo de $\ell_1\vee\dots\vee\ell_k\vee p$ y de $\neg p\vee m_1\vee\dots\vee m_n$. Si $p$ es verdadera en $m$,
$\neg p$ es falsa, así que algún $m_j$ es verdadero; si $p$ es falsa, algún $\ell_i$ es verdadero. En ambos casos
el resolvente es verdadero en $m$. $\blacksquare$

\textbf{(c)} $M(KB\wedge\beta)=M(KB)\cap M(\beta)\subseteq M(KB)\subseteq M(\alpha)$. $\blacksquare$

\textbf{(d)} Resolver $C_1\ni p$ con $C_2\ni\neg p$: $p$ es el único positivo de $C_1$, así que $C_1\setminus\{p\}$ no
tiene positivos; $C_2\setminus\{\neg p\}$ tiene a lo más uno. El resolvente tiene a lo más un literal positivo. $\blacksquare$
\end{solucion}

\begin{solucion}{C16 · Clases de cláusulas y complejidad}
\textbf{(a)} Dos cláusulas no tautológicas son equivalentes sii tienen el mismo conjunto de literales
(una cláusula es falsa exactamente cuando todos sus literales lo son). Con 3 literales sobre 3 símbolos
distintos: $\binom{n}{3}2^3=8\binom n3$; todas las que contienen $p$ y $\neg p$ equivalen a \p{True} (1 clase más). Si se
admiten literales repetidos ($A\vee A\vee B\equiv A\vee B$) se suman $4\binom n2+2n$. En todo caso, $O(n^3)$.

\textbf{(b)} Resolver dos cláusulas de 2 literales da una de $\le2$ literales: la clausura vive en un
conjunto de $O(n^2)$ cláusulas; cada paso agrega una nueva o no agrega nada, hay $O(n^4)$ pares por
revisar: polinomial. Con 3-CNF, resolver dos cláusulas de 3 literales da hasta 4, y así crecen hasta $n$
literales: el número de cláusulas posibles es $3^n$ (cada símbolo positivo, negativo o ausente) y la cota
es exponencial. (3-SAT es NP-completo.)

\textbf{(c)} Cada símbolo entra a la agenda a lo más una vez (tabla de inferidos); al procesarlo, se
decrementa el contador de cada cláusula donde aparece como premisa, una vez por aparición. El trabajo
total es proporcional al número de apariciones de literales: lineal en el tamaño de la KB.

\textbf{(d)} No. Los modelos de un conjunto de cláusulas de Horn son cerrados bajo intersección (AND
bit a bit). $P\vee Q$ tiene los modelos $(1,0)$ y $(0,1)$, cuya intersección $(0,0)$ no lo satisface. Se pierde
la capacidad de expresar información disyuntiva (indefinida); se gana inferencia en tiempo lineal.
\end{solucion}

\begin{solucion}{C17 · Rudolph y Scrooge}
\textbf{(a)} (1) $\forall x\,\p{Niño}(x)\imp\p{Ama}(x,\p{Santa})$; (2) $\forall x\,\forall y\,\p{Ama}(x,\p{Santa})\wedge\p{Reno}(y)\imp\p{Ama}(x,y)$;
(3) $\p{Reno}(\p{Rudolph})\wedge\p{NarizRoja}(\p{Rudolph})$; (4) $\forall x\,\p{NarizRoja}(x)\imp\p{Raro}(x)\vee\p{Payaso}(x)$;
(5) $\forall x\,\p{Reno}(x)\imp\neg\p{Payaso}(x)$; (6) $\forall x\,\p{Raro}(x)\imp\neg\p{Ama}(\p{Scrooge},x)$.
Conclusión: $\neg\p{Niño}(\p{Scrooge})$.

\textbf{(b)} C1 $\neg\p{Niño}(x_1)\vee\p{Ama}(x_1,\p{Santa})$; C2 $\neg\p{Ama}(x_2,\p{Santa})\vee\neg\p{Reno}(y_2)\vee\p{Ama}(x_2,y_2)$;
C3 $\p{Reno}(\p{Rudolph})$; C4 $\p{NarizRoja}(\p{Rudolph})$; C5 $\neg\p{NarizRoja}(x_5)\vee\p{Raro}(x_5)\vee\p{Payaso}(x_5)$;
C6 $\neg\p{Reno}(x_6)\vee\neg\p{Payaso}(x_6)$; C7 $\neg\p{Raro}(x_7)\vee\neg\p{Ama}(\p{Scrooge},x_7)$;
C8 $\p{Niño}(\p{Scrooge})$ (negación de la conclusión).

\textbf{(c)}
\begin{enumerate}[nosep]
\item C8 + C1 $\{x_1/\p{Scrooge}\}$: $\p{Ama}(\p{Scrooge},\p{Santa})$.
\item + C2 $\{x_2/\p{Scrooge}\}$: $\neg\p{Reno}(y_2)\vee\p{Ama}(\p{Scrooge},y_2)$.
\item + C3 $\{y_2/\p{Rudolph}\}$: $\p{Ama}(\p{Scrooge},\p{Rudolph})$.
\item C4 + C5 $\{x_5/\p{Rudolph}\}$: $\p{Raro}(\p{Rudolph})\vee\p{Payaso}(\p{Rudolph})$.
\item + C6 $\{x_6/\p{Rudolph}\}$: $\p{Raro}(\p{Rudolph})\vee\neg\p{Reno}(\p{Rudolph})$; + C3: $\p{Raro}(\p{Rudolph})$.
\item + C7 $\{x_7/\p{Rudolph}\}$: $\neg\p{Ama}(\p{Scrooge},\p{Rudolph})$.
\item Con el paso 3: $\square$. Scrooge no es un niño.
\end{enumerate}
\end{solucion}

\begin{solucion}{C18 · ¿Mató la curiosidad al gato?}
\textbf{LPO.} A: $\forall x\,[\forall y\,\p{Animal}(y)\imp\p{Ama}(x,y)]\imp[\exists y\,\p{Ama}(y,x)]$.
B: $\forall x\,[\exists z\,\p{Animal}(z)\wedge\p{Mata}(x,z)]\imp[\forall y\,\neg\p{Ama}(y,x)]$.
C: $\forall x\,\p{Animal}(x)\imp\p{Ama}(\p{Jack},x)$.
D: $\p{Mata}(\p{Jack},\p{Tuna})\vee\p{Mata}(\p{Curiosidad},\p{Tuna})$.
E: $\p{Gato}(\p{Tuna})$. F: $\forall x\,\p{Gato}(x)\imp\p{Animal}(x)$.
Negación: $\neg\p{Mata}(\p{Curiosidad},\p{Tuna})$.

\textbf{CNF.} A1 $\p{Animal}(F(x))\vee\p{Ama}(G(x),x)$; A2 $\neg\p{Ama}(x,F(x))\vee\p{Ama}(G(x),x)$;
B $\neg\p{Ama}(y,x)\vee\neg\p{Animal}(z)\vee\neg\p{Mata}(x,z)$; C $\neg\p{Animal}(x)\vee\p{Ama}(\p{Jack},x)$;
D, E como arriba; F $\neg\p{Gato}(x)\vee\p{Animal}(x)$; ¬G $\neg\p{Mata}(\p{Curiosidad},\p{Tuna})$.
(Cada cláusula con sus propias variables.)

\textbf{Refutación.}
\begin{enumerate}[nosep]
\item E + F $\{x/\p{Tuna}\}$: $\p{Animal}(\p{Tuna})$.
\item D + ¬G: $\p{Mata}(\p{Jack},\p{Tuna})$.
\item B + (1) $\{z/\p{Tuna}\}$: $\neg\p{Ama}(y,x)\vee\neg\p{Mata}(x,\p{Tuna})$.
\item + (2) $\{x/\p{Jack}\}$: $\neg\p{Ama}(y,\p{Jack})$.
\item A2 + C $\{x/\p{Jack},\ x'/F(\p{Jack})\}$: $\neg\p{Animal}(F(\p{Jack}))\vee\p{Ama}(G(\p{Jack}),\p{Jack})$.
\item + A1 $\{x/\p{Jack}\}$: $\p{Ama}(G(\p{Jack}),\p{Jack})$ (factorizando el literal repetido).
\item + (4) $\{y/G(\p{Jack})\}$: $\square$.
\end{enumerate}
Idea: Jack ama a todos los animales, así que alguien lo ama; pero quien mata un animal no es amado por
nadie; luego Jack no mató a Tuna, y por D fue Curiosidad.
\end{solucion}

\begin{solucion}{C19 · Cuantificadores y la cola de un animal}
\textbf{(a)} Si $M\models\exists y\,\forall x\,\p{Ama}(x,y)$, hay un objeto $b$ con $\p{Ama}(a,b)$ para todo $a$; entonces para
cada $a$ basta tomar $y=b$: $M\models\forall x\,\exists y\,\p{Ama}(x,y)$. Contramodelo del recíproco: dominio $\{1,2\}$,
$\p{Ama}=\{(1,1),(2,2)\}$: todos aman a alguien (a sí mismos) pero nadie es amado por todos.

\textbf{(b)} Premisa: $\forall x\,\p{Perro}(x)\imp\p{Animal}(x)$. Conclusión:
$\forall t\,[\exists x\,\p{Perro}(x)\wedge\p{ColaDe}(t,x)]\imp[\exists y\,\p{Animal}(y)\wedge\p{ColaDe}(t,y)]$.
Negación: $\exists t\,[\exists x\,\p{Perro}(x)\wedge\p{ColaDe}(t,x)]\wedge[\forall y\,\neg\p{Animal}(y)\vee\neg\p{ColaDe}(t,y)]$.
CNF (Skolem $t\to T$, $x\to D$): (1) $\neg\p{Perro}(x)\vee\p{Animal}(x)$; (2) $\p{Perro}(D)$; (3) $\p{ColaDe}(T,D)$;
(4) $\neg\p{Animal}(y)\vee\neg\p{ColaDe}(T,y)$.
Resolución: (2)+(1) $\{x/D\}$: $\p{Animal}(D)$; +(4) $\{y/D\}$: $\neg\p{ColaDe}(T,D)$; +(3): $\square$.
\end{solucion}

\begin{solucion}{C20 · Ensayos (puntos clave)}
\textbf{(a)} La racionalidad está bien definida matemáticamente y es general (incluye conducta sin
inferencia explícita); la prueba de Turing mide parecido humano, que incluye errores humanos y no es un
criterio de diseño ingenieril (los aviones no se diseñan para engañar a las palomas).
\textbf{(b)} Searle sigue un manual para responder en chino sin entenderlo: la sintaxis no basta para la
semántica. Respuesta del sistema: quien entiende no es la persona sino el sistema completo (persona +
manual + papeles); Searle replica que podría memorizar el manual y seguiría sin entender.
\textbf{(c)} No existe algoritmo que decida si un programa arbitrario se detiene; por tanto hay preguntas
bien planteadas que ningún agente computacional puede responder en general; la racionalidad debe ser
\emph{acotada} por los recursos.
\textbf{(d)} Búsqueda sistemática: importa el camino y se requieren garantías (completitud/optimalidad).
CSP: el camino no importa, hay estructura de variables/restricciones explotable (propagación, MRV).
Búsqueda local: espacios enormes u optimización, cuando basta una buena solución y la memoria es limitada.
\end{solucion}

\newpage
% ============================================================
\section{Parte D — Examen simulado: soluciones}

\begin{solucion}{D1 · Jarras (12 pts)}
\textbf{(a)} Estado $(x,y)$: litros en la jarra de 4 y en la de 3. Inicial $(0,0)$. Meta: $x=2\vee y=2$.
Acciones: $\textit{Llenar4}$: $(4,y)$ si $x<4$; $\textit{Llenar3}$: $(x,3)$ si $y<3$; $\textit{Vaciar4}$: $(0,y)$ si $x>0$;
$\textit{Vaciar3}$: $(x,0)$ si $y>0$; $\textit{Verter4}{\to}3$: con $t=\min(x,3-y)$, $(x-t,y+t)$ si $t>0$; $\textit{Verter3}{\to}4$:
con $t=\min(y,4-x)$, $(x+t,y-t)$ si $t>0$. Costo 1.

\textbf{(b)} $5\cdot4=20$ estados; alcanzables \textbf{14}: tras cualquier acción alguna jarra queda llena o vacía,
así que solo son alcanzables los estados con $x\in\{0,4\}$ o $y\in\{0,3\}$ (los 6 estados con $x\in\{1,2,3\}$, $y\in\{1,2\}$ no lo son).

\textbf{(c)} $(0,0)\xrightarrow{\textit{Llenar3}}(0,3)\xrightarrow{3\to4}(3,0)\xrightarrow{\textit{Llenar3}}(3,3)\xrightarrow{3\to4}(4,2)$: \textbf{4 acciones}.
Optimalidad por BFS: nivel 1 $=\{(4,0),(0,3)\}$; nivel 2 $=\{(1,3),(4,3),(3,0)\}$; nivel 3 $=\{(1,0),(3,3)\}$.
Ningún estado de nivel $\le3$ contiene 2 litros, así que no hay solución más corta.
\end{solucion}

\begin{solucion}{D2 · A* y consistencia (15 pts)}
\textbf{(a)} $h^*$: S=7, A=6, B=4, C=3, G=0. Admisible ($6\le7$, $5\le6$, $2\le4$, $3\le3$). Arcos:
S$\to$A $6\le1+5$; S$\to$B $6\le4+2$; \textbf{A$\to$B: $5>2+2$} \texttimes{}; A$\to$C $5\le5+3$; B$\to$C $2\le1+3$;
B$\to$G $2\le6$; C$\to$G $3\le3$. \textbf{Inconsistente} por A$\to$B.

\textbf{(b)}
\begin{center}\small
\begin{tabular}{clll}
\toprule
Paso & Extrae ($g$,$f$) & Frontera & Comentario\\
\midrule
1 & S (0,6) & A(6), B(6) & \\
2 & A (1,6) & B(5), C(9) & B mejora a $g=3$\\
3 & B (3,5) & C(7), G(9) & C mejora a $g=4$\\
4 & C (4,7) & G(7) & G mejora a $g=7$\\
5 & G (7,7) & — & costo \textbf{7} (óptimo)\\
\bottomrule
\end{tabular}
\end{center}

\textbf{(c)} Con B primero: S; B$(g=4,f=6)$ se cierra con $g$ subóptimo, genera C$(5,8)$ y G$(10,10)$; A$(1,6)$
genera B con $g=3$, pero B ya está cerrado y se descarta; C vía A ($g=6$) no mejora; C$(5,8)$ genera G$(8,8)$;
extrae G: \textbf{costo 8}, subóptimo. La inconsistencia hace que $f$ baje a lo largo de S–A–B
($6\to6\to5$), así que B puede cerrarse antes de hallar su mejor camino; sin reapertura, el error se propaga.

\textbf{(d)} Ver C4(a).
\end{solucion}

\begin{solucion}{D3 · CSP (13 pts)}
\textbf{(a)} Ciclo de 4: $A$–$B$–$C$–$D$–$A$ (aristas $A<B$, $B<C$, $C\ne D$, $D>A$).

\textbf{(b)} Una ejecución de AC-3 (cola inicial en orden alfabético de arcos):
$A\to B$ quita 3 ($A=\{1,2\}$); $B\to A$ quita 1 ($B=\{2,3\}$); $B\to C$ quita 3 ($B=\{2\}$);
$C\to B$ quita 1, 2 ($C=\{3\}$); $D\to A$ quita 1 ($D=\{2,3\}$); $D\to C$ quita 3 ($D=\{2\}$);
reencolado $A\to B$ quita 2 ($A=\{1\}$). Final: $A=1$, $B=2$, $C=3$, $D=2$. Todos unitarios: es
solución, \textbf{no hace falta buscar} (verificación: $1<2<3$, $2\ne3$, $2>1$).

\textbf{(c)} Todos los dominios miden 3: MRV empata; todos tienen grado 2: empate otra vez; se toma $A$ (alfabético).

\textbf{(d)} Tras $A=1$, FC solo poda los vecinos de $A$: $B=\{2,3\}$, $D=\{2,3\}$, $C=\{1,2,3\}$. AC-3 propaga en
cadena: $C>B\ge2\Rightarrow C=\{3\}$; $B<3\Rightarrow B=\{2\}$; $D\ne3\Rightarrow D=\{2\}$: resuelve el problema.
\end{solucion}

\begin{solucion}{D4 · Alfa-beta (15 pts)}
\textbf{(a)} MIN de nivel 3: $111=2$, $112=1$, $121=4$, $122=3$, $211=2$, $212=1$, $221=5$, $222=3$.
MAX: $11=2$, $12=4$, $21=2$, $22=5$. MIN: $1=2$, $2=2$. Raíz \textbf{2}.

\textbf{(b)}
\begin{itemize}[nosep]
\item $R(-\infty,\infty)\to1(-\infty,\infty)\to11(-\infty,\infty)\to111(-\infty,\infty)$: 2, 5 $\Rightarrow2$. En 11: $\alpha=2$.
\item $112(2,\infty)$: 8, 1 $\Rightarrow v=1\le\alpha$ (corte al final). $11=2$. En 1: $\beta=2$.
\item $12(-\infty,2)\to121(-\infty,2)$: 4, 6 $\Rightarrow4$. En 12: $v=4\ge\beta=2$ $\Rightarrow$ \textbf{se poda 122} (hojas 7, 3). $1=2$.
\item En $R$: $\alpha=2$. $2(2,\infty)\to21(2,\infty)\to211(2,\infty)$: 9, 2 $\Rightarrow v=2\le\alpha$: devuelve 2.
\item $212(2,\infty)$: 1 $\le\alpha$ $\Rightarrow$ \textbf{se poda la hoja 6}. $21=2$.
\item En 2: $v=2\le\alpha=2$ $\Rightarrow$ \textbf{se poda 22} (hojas 5, 8, 3, 4). $R=2$.
\end{itemize}
Hojas no evaluadas: 7, 3 (de 122), 6 (de 212), 5, 8, 3, 4 (de 22): \textbf{7 de 16}; se evalúan 9.
(Con la convención de podar solo con desigualdad estricta, el subárbol 22 sí se exploraría.)

\textbf{(c)} $b^{\lceil m/2\rceil}+b^{\lfloor m/2\rfloor}-1=4+4-1=\mathbf{7}$ hojas.
\end{solucion}

\begin{solucion}{D5 · Búsqueda local (10 pts)}
\textbf{(a)} $\Delta E=-3$: $T=5$: $e^{-0.6}\approx0.549$; $T=1$: $e^{-3}\approx0.050$. Con $T\to0$ la probabilidad
tiende a 0: se vuelve hill climbing.

\textbf{(b)} \texttt{110|10010} $\times$ \texttt{011|01101} $\Rightarrow$ \texttt{11001101} (5 unos) y \texttt{01110010}
(4 unos); los padres tienen 4 y 5. $P(\text{sin mutación})=(7/8)^8\approx0.344$.

\textbf{(c)} Tabú se mueve al mejor vecino no tabú aunque empeore y no puede regresar de inmediato, así
que sale de la cuenca; hill climbing se detiene si ningún vecino mejora.
\end{solucion}

\begin{solucion}{D6 · Lógica proposicional (15 pts)}
\textbf{(a)} Si el consecuente fuera falso ($Q=S=F$) con el antecedente verdadero: $P\imp Q$ fuerza $P=F$,
$R\imp S$ fuerza $R=F$, y entonces $P\vee R$ es falso. Contradicción: no hay asignación que la falsifique; es válida.

\textbf{(b)} (1) $\neg P\vee Q$, (2) $\neg R\vee S$, (3) $P\vee R$, (4) $\neg Q$, (5) $\neg S$.
(1)+(4): $\neg P$; (2)+(5): $\neg R$; (3)+$\neg P$: $R$; $R+\neg R$: $\square$.

\textbf{(c)} No: $P\vee R$ tiene dos literales positivos y la KB no es equivalente a ninguna KB de Horn
(sus modelos $\{P,Q\}$ y $\{R,S\}$ tienen intersección vacía, que no satisface $P\vee R$). Modus ponens no
puede concluir $Q\vee S$ porque no conoce $P$ ni $R$ por separado: se necesita razonar por casos
(resolución).

\textbf{(d)} $3\cdot2^2=\mathbf{12}$.
\end{solucion}

\begin{solucion}{D7 · Lógica de primer orden (20 pts)}
\textbf{(a)} $\forall x\,\p{SabeLeer}(x)\imp\p{Letrado}(x)$; $\forall x\,\p{Delfin}(x)\imp\neg\p{Letrado}(x)$;
$\exists x\,\p{Delfin}(x)\wedge\p{Inteligente}(x)$. Conclusión: $\exists x\,\p{Inteligente}(x)\wedge\neg\p{SabeLeer}(x)$.

\textbf{(b)} C1 $\neg\p{SabeLeer}(x)\vee\p{Letrado}(x)$; C2 $\neg\p{Delfin}(y)\vee\neg\p{Letrado}(y)$;
C3 $\p{Delfin}(F_1)$; C4 $\p{Inteligente}(F_1)$ ($F_1$ constante de Skolem del $\exists$). Negación de la conclusión:
$\forall z\,\neg\p{Inteligente}(z)\vee\p{SabeLeer}(z)$ $\Rightarrow$ C5 $\neg\p{Inteligente}(z)\vee\p{SabeLeer}(z)$ (aquí no hay Skolem).

\textbf{(c)} C4+C5 $\{z/F_1\}$: $\p{SabeLeer}(F_1)$; +C1 $\{x/F_1\}$: $\p{Letrado}(F_1)$; +C2 $\{y/F_1\}$: $\neg\p{Delfin}(F_1)$;
+C3: $\square$.

\textbf{(d)} $\{x/\p{Padre}(y),\ z/\p{Madre}(\p{Padre}(y))\}$. $P(x,x)$ y $P(F(y),y)$: $x/F(y)$ y luego $F(y)$ vs.\ $y$:
la verificación de ocurrencia falla.
\end{solucion}

\begin{nota}
\textbf{Rúbrica sugerida para autocalificarte.} En cada inciso: 40\% planteamiento/definiciones,
40\% procedimiento, 20\% resultado. Si obtuviste menos de 70 puntos, repite los ejercicios de las
Partes B y C del tema correspondiente antes del examen.
\end{nota}


\end{document}
